\documentclass{article}
\usepackage{arxiv}

% \documentclass[a4paper,14pt]{extarticle}
% \usepackage[a4paper,top=2.0cm, bottom=1.5cm, left=3.0cm, right=1.5cm]{geometry}

\usepackage[utf8x]{inputenc}
\usepackage[english, russian]{babel}
% \usepackage[T2A, T1]{fontenc}
\usepackage{url}
\usepackage{booktabs}
\usepackage{amsfonts}
\usepackage{nicefrac}
\usepackage{microtype}
\usepackage{lipsum}
\usepackage{graphicx}
\usepackage{natbib}
\usepackage{doi}
\usepackage{amsmath,amsfonts,amssymb,amsthm,mathtools}
\usepackage{caption}
\usepackage{float}
\usepackage{wrapfig}
\usepackage{subcaption}
\usepackage{xcolor}
\usepackage{amsthm}
\usepackage{cleveref}
\usepackage{algorithm}
\usepackage{algpseudocode}
\usepackage{hyperref}


% \linespread{1.5}

%\usepackage{subfigure}
\DeclareMathOperator*{\argmax}{\arg\!\max}
\DeclareMathOperator*{\argmin}{\arg\!\min}
\DeclareMathOperator{\tr}{tr}
\DeclareMathOperator{\rank}{rank}
\DeclareMathOperator{\sr}{sr}



\title{Stable Rank Based Analysis of Low-Rank Adaptation: Tighter Expressivity and Generalization Bounds}

\author{Nikita Okhotnikov\\
	MIPT \\
\and
	Andrey Grabovoy\\
	MIPT \\
}
\date{}

% \renewcommand{\shorttitle}{Stable Rank Analysis of LoRA}

\hypersetup{
pdftitle={Stable Rank Analysis of Low-Rank Adaptation},
pdfsubject={Machine Learning, Parameter-Efficient Fine-Tuning},
pdfauthor={Nikita Okhotnikov},
}

\newtheorem{theorem}{Theorem}
\newtheorem{lemma}{Lemma}
\newtheorem{corollary}{Corollary}
\newtheorem{definition}{Definition}
\newtheorem{remark}{Remark}


\begin{document}
    \maketitle
    \begin{abstract}
        % ==============================================================================
        % IN SHORT:
        % - LoRA is the standard PEFT method, but its theoretical bounds rely heavily on algebraic rank.
        % - Empirically, trained LoRA adapters exhibit a stable rank (SR) significantly lower than their nominal rank.
        % - We propose stable rank as a mathematically robust measure for bounding LoRA's capacity.
        % - We derive fundamental limits on matrix approximation under SR constraints, extending these to expressivity bounds for FNNs and Transformers.
        % - Using a PAC-Bayesian framework, we prove a tighter generalization bound scaling as O(\sqrt{sr(\Delta W)/N}), an improvement over prior O(\sqrt{rank/N}) bounds.
        % - We empirically validate our theoretical findings, confirming LoRA's implicit bias toward low stable rank and demonstrating that the generalization gap practically correlates with \sqrt{sr}.
        % ==============================================================================
        
        Low-rank adaptation (LoRA) has become the de facto standard for parameter-efficient fine-tuning of large language models. 
        Existing theoretical bounds on expressivity and generalization of LoRA are formulated in terms of algebraic rank, which fails to capture the empirical observation that trained matrices often exhibit stable rank significantly lower than their nominal rank. 
        In this paper, we propose stable rank as a more appropriate theoretical measure for analyzing LoRA's properties. 
        We derive fundamental limits on single matrix approximation under stable rank constraint, extend these to expressivity bounds for fully connected networks and transformers. 
        Additionally, we prove a tighter generalization bound proportional to $\sqrt{\sr(\Delta W)}$ rather than previous best $\sqrt{\rank(\Delta W)}$ from the literature, using a PAC-Bayesian framework. 
        We manage to validate our theoretical findings and demonstrate via controlled experiments on synthetic data that $\mathcal{O}(\sqrt{\sr(\Delta W)})$ is more practical estimate for generalization gap. Our experiments additionally confirms that lora exhibits some implicit bias toward lower stable rank solutions even when the rank of the adapter becomes as large as the width of the frozen layer.
    \end{abstract}

    \tableofcontents
    \vspace{1cm}


    %%%%%%%%%%%%%%%%%%%%%%%%%%%%%%%%%%%%%%%%%%%%%%%%%%%%%%%%%%%%%%%%%%%%%%%%%%%%%%%%
    \section{Introduction}
    % ==============================================================================
    % IN SHORT:
    % - Expanding LLM sizes make full fine-tuning computationally infeasible.
    % - LoRA allows efficient adaptation, yet theoretical bounds often misalign with empirical reality because they rely on algebraic rank.
    % - Algebraic rank overestimates effective capacity since optimization naturally skews energy toward a few dominant directions.
    % - Stable rank (SR), which accounts for energy distribution across singular values, aligns better with observed LoRA behavior.
    % - We introduce theorems proving that the true expressivity of LoRA is bounded by its SR constraint.
    % - We present a novel PAC-Bayesian generalization bound that swaps nominal rank for SR, tightening the estimation of overfitting.
    % - The implicit optimization bias causing this low SR phenomenon is acknowledged but reserved for future work.
    % ==============================================================================
    

    Recent advances in large language models (LLMs) and multimodal models (MLLMs) induce a growing number of various downstream tasks attempted to be solved with a single pretrained model. 
    Despite the fact that the performance of such models has vastly improved over recent years, there are still tasks that cannot be handled in a zero-shot setting and others that heavily benefit from additional fine-tuning. 
    However, the size of LLMs with cutting-edge performance keeps growing, and even the best open-sourced models nowadays contain hundreds of billions of parameters \cite{deepseekai2025deepseekr1incentivizingreasoningcapability,kimiteam2025kimik2openagentic,yang2025qwen3technicalreport,grattafiori2024llama3herdmodels}. 
    Thus, full fine-tuning of such models is computationally infeasible in many cases, making parameter-efficient fine-tuning (PEFT) methods strictly required. 
    Among these, low-rank adaptation (LoRA) \cite{hu2021loralowrankadaptationlarge} has emerged as the most widely adopted technique, enabling the efficient adaptation of billion-parameter models by freezing pretrained weights and injecting small low-rank  rank trainable matrices. 
    Since the release of the original paper, the approach has been actively studied \cite{biderman2024lora, bernoulli2025lora, zeng2024expressivepowerlowrankadaptation, kratsios2025sharpgeneralizationboundsfoundation,jang2024why, malinovsky2024randomized}, yet it still lacks a comprehensive theoretical foundation explaining its expressivity, generalization, expected optimization behavior and some emprically observable effects.
    
    Most previous theoretical results on LoRA are based on the nominal algebraic rank of the adapters as the primary bounding parameter. 
    However, such rank-based approaches diverge from empirical observations: trained LoRA adapters often produce very specific matrices that exhibit a stable rank close to 1 \cite{lion2025polar} on various downstream tasks. 
    The stable rank of a matrix $A \in \mathbb{R}^{m \times n}$ is defined as:
    \begin{equation}
        \sr(A) = \frac{\|A\|_F^2}{\|A\|_{op}^2} = \frac{\sum_{i=1}^{\min(m,n)} \sigma_i^2(A)}{\sigma_1^2(A)},
    \end{equation}
    where $\sigma_i(A)$ are the singular values in descending order. 
    This quantity, introduced by Rudelson and Vershynin \cite{rudelson2007sampling}, provides a continuous measure of the ``effective dimensionality'' of a matrix. 
    Unlike algebraic rank, stable rank is robust to small perturbations and captures the distribution of energy across singular values. 
    The discrepancy between theoretical bounds based on algebraic rank and empirical observations comes from a fundamental mismatch in how we estimate LoRA's capacity. 
    While algebraic rank provides an upper bound on the number of independent directions a matrix can span, it fails to capture how the gradient-based optimization procedure actually works with these parameters' theoretical capacity. 

    In this work, we claim that stable rank is a more mathematically appropriate measure for the theoretical analysis of low-rank adapters, and we formally study the expressivity and generalization properties of the trained adapter under this lens. 
    We additionally confirm the implicit bias of low stable rank in both controlled synthetic experiments and real-world tasks. 
    Unlike previous works, we operate strictly under a stable rank constraint and focus on the dependency of LoRA's properties on the resulting stable rank. 
    We prove a fundamental lemma on the irreducible approximation error of a single matrix by an estimate of limited stable rank, then providing an actual tight lower bound and the way to construct the matrix that reaches that bound similar to the Eckart-Young-Mirsky theorem \cite{EckartYoung1936}. Using this lemma, we derive a novel, more realistic expressivity upper bound that scales with stable rank rather than algebraic rank for both multi-layer feed-forward networks (FNNs) and slightly simplified transformer \cite{vaswani2023attentionneed} architecture. 
    Our results incorporate this implicit optimization bias into previous findings on expressivity limits. 

    Additionally, we demonstrate that the degree of overfitting -- or the size of the generalization gap -- is better characterized as a function of stable rank rather than nominal rank.
    Utilizing a PAC-Bayesian framework \cite{mcallester1999pac, mcallester2003simplified}, we provide a rigorous proof for a generalization gap upper bound that scales as $\mathcal{O}(\sqrt{\sr(\Delta W)/N})$. 
    This significantly tightens the previously known best bound of $\mathcal{O}(\sqrt{\rank(\Delta W)/N})$ \cite{kratsios2025sharpgeneralizationboundsfoundation}.
    
    The root cause of the observed low stable rank implicit bias is beyond the scope of this paper, as is an exhaustive analysis of the optimization trajectory and procedure itself. 
    This remains an active subject for future research.


    %%%%%%%%%%%%%%%%%%%%%%%%%%%%%%%%%%%%%%%%%%%%%%%%%%%%%%%%%%%%%%%%%%%%%%%%%%%%%%%%
    \section{Theoretical Background}
    % ==============================================================================
    % IN SHORT:
    % - Reviews PEFT methods transitioning from adapters to prompt tuning to LoRA.
    % - Highlights existing LoRA theory (Zeng & Lee on expressivity, Kratsios on generalization bounds), noting their reliance on nominal rank.
    % - Discusses the rise of Stable Rank in ML theory (Tropp, Bartlett) and recent findings like PoLAR mapping its necessity for LoRA.
    % - Explains the PAC-Bayesian framework which natively connects KL divergence of weight posteriors to spectral/Frobenius norms.
    % ==============================================================================


    %%%%%%%%%%%%%%%%%%%%%%%%%%%%%%%%%%%%%%%%%%%%%%%%%%%%%%%%%%%%%%%%%%%%%%%%%%%%%%%%
    \subsection{Parameter-Efficient Fine-tuning}

    Parameter-Efficient Fine-tuning (PEFT) has emerged as a crucial paradigm in modern deep learning, designed to adapt massive pre-trained models to downstream tasks while updating only a minimal subset of parameters. 
    The landscape of PEFT methods has diversified significantly over time, with each method offering different trade-offs between parameter efficiency, training speed, inference latency, and task performance.

    Before the dominance of LoRA, the primary methods broadly fell into architectural additions and input modifications. 
    Houlsby et al. \cite{houlsby2019parameter} introduced bottleneck adapters that insert small MLP layers between existing transformer blocks, while Rebuffi et al. \cite{rebuffi2017learning} explored learning multiple adapters for multi-task scenarios. 
    These methods achieve good parameter efficiency but fundamentally alter the forward pass, adding unavoidable inference latency due to the sequential nature of the additional layers. 
    Alternatively, Li et al. \cite{li2021prefix} proposed prefix-tuning that optimizes continuous prefixes of input sequences, while Lester et al. \cite{lester2021power} introduced prompt tuning with learned soft prompts. Qin et al. \cite{qin2022towards} provided a unified framework analyzing the relationship between these prompt-based approaches. 
    While these methods modify inputs rather than weights, avoiding most of theinference latency, they consume valuable context window space and present different theoretical limitations.

    LoRA \cite{hu2021loralowrankadaptationlarge} established a new standard by hypothesizing that changes in model weights during fine-tuning have a low "intrinsic dimension." 
    By freezing the pre-trained weights $W_0 \in \mathbb{R}^{d \times k}$ and injecting trainable rank decomposition matrices $A \in \mathbb{R}^{r \times k}$ and $B \in \mathbb{R}^{d \times r}$ such that the update is $\Delta W = BA$, LoRA avoids both additional inference latency (as weights can be merged into pretrained modelsafter training) and context window reductions. Additionally, as some other methods, it allows learning and inferencing multiple adapters for different tasks without re-hosting and resaving the entire pretrained model.

    Because of its structural simplicity and lack of inference overhead, LoRA has become the de facto standard PEFT method. Comprehensive comparative studies by Karimi et al. \cite{karimi2021compelling} and He et al. \cite{he2021effectiveness} empirically verified its strong performance across varying tasks and model scales. Extensive surveys by Ding et al. \cite{ding2023parameter}, Han et al. \cite{han2024survey}, and Xu et al. \cite{xu2023parameter} further document its ubiquity across the PEFT varieties. 
    
    Following its success, numerous structural and optimization-based variants have emerged to address specific limitations of the vanilla formulation. Zhang et al. \cite{zhang2023adaptive} developed AdaLoRA to dynamically allocate rank budgets across different layers based on importance scores, while Valipour et al. \cite{valipour2023dylora} introduced DyLoRA for dynamic, search-free rank selection during training. To improve the learning dynamics, Liu et al. \cite{liu2024lora} proposed LoRA+, which applies asymmetric learning rates to the $A$ and $B$ matrices to overcome initialization-induced optimization inefficiencies. Kalajdzievski \cite{kalajdzievski2023rankstabilizationscalingfactor} introduced Rank-Stabilized LoRA (RS-LoRA), which modifies the scaling factor from $1/r$ to $1/\sqrt{r}$ to mathematically stabilize the learning trajectory and prevent gradient collapse at high ranks. Liu et al. \cite{liu2024doraweightdecomposedlowrankadaptation} proposed Weight-Decomposed Low-Rank Adaptation (DoRA), which decouples the pretrained weights into magnitude and directional components, applying the standard LoRA update strictly to the directional matrix to closer mimic full fine-tuning behavior. Meng et al. \cite{meng2025pissaprincipalsingularvalues} introduced PiSSA, which initializes the $A$ and $B$ matrices using the Principal Singular Values and Vectors of the base model weights instead of random noise, allowing the adapter to immediately optimize the most significant spectral components. 

    Despite these seeming diversity -- whether a method alter the initialization (PiSSA), decouple the weight vectors (DoRA), stabilize the scaling (RS-LoRA), or adjust the learning rates (LoRA+) -- all of them ultimately preserve the core low-rank multiplicative structure $\Delta W = BA$ applied additively to the base weights. Because our mathematical framework operates directly on these matrices and the stable rank of the resulting perturbation matrix $\Delta W$, rather than relying on assumptions about the specific optimizer state or initialization scheme, the theoretical principles and bounds presented in this paper remain fully applicable to these variants.



    %%%%%%%%%%%%%%%%%%%%%%%%%%%%%%%%%%%%%%%%%%%%%%%%%%%%%%%%%%%%%%%%%%%%%%%%%%%%%%%%   
    \subsection{LoRA Theoretical Background}
    
    The theoretical understanding of LoRA has developed along several complementary directions, though most share a common limitation in their reliance on the algebraic rank parameter $r$. 

    Regarding expressive power, Zeng and Lee \cite{zeng2024expressivepowerlowrankadaptation} established a foundational existence proof showing that rank-$r$ LoRA adapters can represent any target perturbation of a given width and depth when certain rank threshold is met. However, this existence result operates purely in the algebraic capacity space, it provides no insight into the actual penalty incurred if gradient-based optimization converges to matrices with skewed singular value distributions or any insights onto the form of this disctributions in the <<optimal>> solution. 
    
    From a generalization perspective, Kratsios et al. \cite{kratsios2025sharpgeneralizationboundsfoundation} derived sample complexity and generalization bounds specifically for asymmetric low-rank adapters. Rather than relying on guarantees averaged over multiple random seeds, they studied the concentration of the generalization gap for a single fine-tuning run, yielding sharp, high-probability upper bounds that scale proportionally to $\mathcal{O}(\sqrt{r/N})$ along with a matching theoretical lower bound. Wu et al. \cite{wu2024generalization} also analyzed LoRA generalization but through the lens of Neural Tangent Kernels (NTK). By linearizing the non-linear network dynamics in the wide-model limit, they demonstrated that choosing an adapter rank $r \gtrsim \sqrt{N}$ guarantees the elimination of spurious local minima while maintaining generalization performance comparable to full fine-tuning in underdetermined regimes. 

    Additionally, the optimization dynamics of factored weight matrices have been heavily studied. Wang et al. \cite{wang2024mechanism} explored the implicit bias of LoRA optimization, revealing that the non-convex $A \times B$ factorization intrinsically steers gradient descent toward well-regularized solutions, which acts as a mechanism to mitigate catastrophic forgetting without requiring explicit penalty terms. Tian et al. \cite{tian2024optimization} analyzed convergence properties under different initialization schemes, theoretically demonstrating that standard random initialization often traps LoRA in suboptimal low-rank subspaces; they showed that reaching the optimal minimizers often strictly requires data-driven or ``preheated'' initializations. Zhu et al. \cite{zhu2024asymmetry} explored the theoretical implications of asymmetric structures in the low-rank factors, proving that the down-projection and up-projection matrices exhibit fundamentally different spectral dynamics and learning speeds during optimization. Finally, Liu et al. \cite{liu2023convergence} provided robust convergence guarantees for LoRA under overparameterized settings, leveraging Polyak-Lojasiewicz (PL) conditions to theoretically guarantee that global minima remain accessible via gradient descent, provided the base network is sufficiently wide.

    Despite that variety of approaches tested in previous works, a critical element is missing: these results largely fail to account for the actual spectral structure and energy distribution of the learned adapters. Algebraic rank assumes a maximum entropy scenario where all $r$ components contribute equally or at least expected to be able to contribute equally. Recently, Lion et al. \cite{lion2025polar} identified this blind spot in their paper introducing PoLAR (Polar-Decomposed Low-Rank Adapter Representation). They demonstrated empirically that standard LoRA suffers from a "directional diversity collapse," learning updates with stable rank well below the regular algebraic rank of the adapters, often even close to 1. This phenomenon severely degrades the effective utilization of the allocated parameter subspace. Inspired by these practical observations, our work bridges this gap by replacing nominal rank with stable rank directly in the theoretical bounds of expressivity and generalization.


    %%%%%%%%%%%%%%%%%%%%%%%%%%%%%%%%%%%%%%%%%%%%%%%%%%%%%%%%%%%%%%%%%%%%%%%%%%%%%%%%   
    \subsection{Stable Rank in Machine Learning}
    
    Stable rank has historically emerged as a powerful, continuous tool for analyzing matrix properties in machine learning and random matrix theory, offering distinct advantages over algebraic rank. 

    Introduced formally by Rudelson and Vershynin \cite{rudelson2007sampling} in the context of sampling from large matrices, stable rank serves as a continuous measure of effective dimensionality. Tropp \cite{tropp2012user, tropp2015introduction} expanded heavily on its properties, showing that matrix concentration inequalities in random matrix theory are far more tightly bounded by stable rank than by algebraic dimensions. In the neural network theory, Bartlett et al. \cite{bartlett2017spectrally} and Neyshabur et al. \cite{neyshabur2015norm} utilized spectral norms and matrix energy distributions to derive generalization bounds for deep networks, proving that the Lipschitz continuity and spectral complexity of a network better predict generalization than parameter counting.

    More recently, stable rank has been also applied to modern transformer architectures. As mentioned, the PoLAR authors \cite{lion2025polar} documented the tendency of LoRA to collapse into an effective stable rank close to 1. Concurrently, Saada et al. \cite{saada2025mindgapspectralanalysis} utilized Random Matrix Theory to analyze rank collapse in attention layers, showing that standard softmax attention inherently develops a widening spectral gap between the largest eigenvalue and the rest, leading to rank collapse in width. Furthermore, spectral norm regularizations have been effectively utilized by Miyato et al. \cite{miyato2018spectral} and Yoshida et al. \cite{yoshida2017spectral} to improve generalization by artificially constraining the same metrics that stable rank relies on. These works together suggest that stable rank is not just an alternative metric, but the governing mathematical variable dictating the actual complexity and robustness of the networks, thus it is actually useful for adapters analysis.



    %%%%%%%%%%%%%%%%%%%%%%%%%%%%%%%%%%%%%%%%%%%%%%%%%%%%%%%%%%%%%%%%%%%%%%%%%%%%%%%%  
    \subsection{PAC-Bayesian Framework}
    
    To properly integrate stable rank into a generalization bound, we require a framework that natively incorporates continuous, energy-based matrix norms rather than combinatorial parameter counting (like VC-dimension). The PAC-Bayesian framework provides exactly this principled approach to deriving data-dependent generalization bounds.

    The foundational PAC-Bayes theorem, established by McAllester \cite{mcallester1999pac, mcallester2003simplified} and extended to neural networks by Shawe-Taylor and Williamson \cite{shawe1997pac}, bounds the expected generalization gap of a randomized predictor. Let $\mathcal{P}$ be a prior distribution over hypotheses chosen before seeing the training data, and $\mathcal{Q}$ be a posterior distribution learned after observing data. With probability at least $1-\delta$ over a dataset of size $N$, the true expected risk $L_0$ is bounded by the empirical risk $\widehat{L}$ plus a complexity penalty:
    \begin{equation}
        \mathbb{E}_{w \sim \mathcal{Q}}[L_{0}(w)] \leqslant \mathbb{E}_{w \sim \mathcal{Q}}[\widehat{L}(w)] + \sqrt{\frac{KL(\mathcal{Q} || \mathcal{P}) + \ln(2N/\delta)}{2N}}
    \end{equation}
    where $KL(\mathcal{Q} || \mathcal{P})$ is the Kullback-Leibler divergence between the posterior and the prior. 

    This framework is exceptionally well-suited for fine-tuning scenarios, including Low-rank adaptation. If we set the prior $\mathcal{P}$ to be an isotropic Gaussian distribution centered on the pre-trained frozen weights $W_0$, and the posterior $\mathcal{Q}$ to be a distribution centered on the adapted weights $W_0 + \Delta W$, the KL divergence cleanly reduces to the scaled Frobenius norm of the perturbation: $KL \propto \|\Delta W\|_F^2$. 
    
    Building on this, Neyshabur et al. \cite{neyshabur2017pac} and Bartlett et al. \cite{bartlett2017spectrally} developed spectrally-normalized margin bounds for deep networks by combining the PAC-Bayes theorem with Lipschitz continuity across layers. 
    
    More recently, Kratsios et al. \cite{kratsios2025sharpgeneralizationboundsfoundation} established the previous best sample complexity bound of $\tilde{\mathcal{O}}(\sqrt{r/N})$ for asymmetric randomized low-rank adapters. However, rather than utilizing a PAC-Bayesian approach, their analysis relies on constructive approximation theory and empirical Rademacher complexity. Specifically, their method bounds the local Lipschitz regularity of the parameter-to-network map, uses this to estimate the covering numbers of the randomized LoRA function class, and finally applies Dudley's entropy integral to bound the generalization gap uniformly over the entire hypothesis space. While their approach yields a mathematically sharp sample complexity, we cannot follow this same path. As a result of their framework, the topological dimension of the covering space is intrinsically tied to the ambient parameter count. This lead to an unbreakable proportionality to the algebraic rank, resulting in the $\mathcal{O}(\sqrt{r/N})$ scaling. If we were to introduce a stable rank constraint within their volumetric covering-number approach, it would not effectively reduce the topological dimension of the parameter space, making it theoretically impossible to decouple the final generalization bound from $\sqrt{r}$. In contrast, we adopt the PAC-Bayesian framework because its complexity measure--the Kullback-Leibler divergence between a data-independent prior and the learned posterior--naturally captures the spectral and norm-based properties of the network weights. This allows us to smoothly integrate the stable rank constraint directly into the KL-divergence term, successfully bypassing the ambient algebraic rank entirely to yield a tighter generalization bound that scales with the stable rank rather than $r$.



    %%%%%%%%%%%%%%%%%%%%%%%%%%%%%%%%%%%%%%%%%%%%%%%%%%%%%%%%%%%%%%%%%%%%%%%%%%%%%%%%
    \section{Mathematical Framework}
    % ==============================================================================
    % IN SHORT:
    % - Formalizes notation (Frobenius norm, operator norm, Stable Rank).
    % - Derives a lower bound (Lemma 1) proving irreducible error when approximating a high-SR matrix with a low-SR one.
    % - Derives a constructive upper bound (Lemma 2) showing how to build an approximation that obeys a strict SR limit.
    % - Applies these matrix-level bounds to FNNs (Theorem 1), showing layer-wise error amplification of the SR constraint.
    % - Extends the expressivity bound to Transformer architectures (Theorem 2), managing the complex dual-path errors of MHA and FFN.
    % - Concludes with Generalization (Theorem 3), executing a PAC-Bayesian proof to establish the O(\sqrt{sr/N}) bound.
    % ==============================================================================


    %%%%%%%%%%%%%%%%%%%%%%%%%%%%%%%%%%%%%%%%%%%%%%%%%%%%%%%%%%%%%%%%%%%%%%%%%%%%%%%%
    \subsection{Preliminaries}
    
    We use standard matrix norm notation: $\|A\|_F = \sqrt{\sum_{i,j} A_{ij}^2} = \sqrt{\sum_{i} \sigma_i^2(A)}$ denotes the Frobenius norm, and $\|A\|_{op} = \sigma_1(A)$ denotes the spectral (operator) norm. For a matrix $A \in \mathbb{R}^{m \times n}$ with singular values $\sigma_1 \geq \sigma_2 \geq \cdots \geq \sigma_{\min(m,n)} \geq 0$, the stable rank is defined as $\sr(A) = \|A\|_F^2 / \|A\|_{op}^2$. Note that $1 \leq \sr(A) \leq \rank(A) \leq \min(m, n)$. 

    To provide intuition, if $A$ is a dense matrix sampled from a standard normal distribution $\mathcal{N}(0, 1)$, its stable rank is typically proportional to its algebraic dimension (specifically $\sr(A) \approx \min(m,n)$ for large matrices). Thus, observing a stable rank much closer to 1 in trained weights represents a significant, non-trivial collapse of energy into a few dominant directions, rather than a default mathematical property of all matrices.
    
    Throughout this work for a linear layer with weight matrix $W$, we consider a LoRA adapter parameterized as $\Delta W = BA$ where $B \in \mathbb{R}^{m \times r}$, $A \in \mathbb{R}^{r \times n}$, and $r$ is the nominal algebraic rank.



    %%%%%%%%%%%%%%%%%%%%%%%%%%%%%%%%%%%%%%%%%%%%%%%%%%%%%%%%%%%%%%%%%%%%%%%%%%%%%%%%
    \subsection{Single Matrix Approximation}
    
    We begin by analyzing the fundamental problem of approximating a target matrix $A$ with an estimate $\hat{A}$ under stable rank constraints. This setting isolates the core empirical issue that LoRA faces: even if the nominal rank $r$ is large, the optimization process tends to produce matrices with constrained stable rank, limiting their ability to fit arbitrary targets.
    
    The classical Eckart-Young-Mirsky theorem \cite{EckartYoung1936} provides the foundation for understanding low-rank matrix approximation. It states that for any matrix $A$ and integer $k$, the optimal rank-$k$ approximation in the Frobenius norm is obtained by truncating the Singular Value Decomposition (SVD) to the top $k$ singular values: $\hat{A}_k = \sum_{i=1}^k \sigma_i u_i v_i^T$. However, our setting introduces a crucial distinction. Instead of merely constraining the number of non-zero singular values (algebraic rank), we constrain their distribution through the stable rank $\sr(\hat{A}) = \frac{\sum \sigma_i^2}{\sigma_1^2} \leq S$. Under an algebraic rank constraint, the optimal solution zeros out tail singular values. Under a stable rank constraint, one must actively concentrate the energy, either by shrinking the tail or exploding the top singular value. \\

    \begin{lemma} \textbf{(Interpretable SR-Constrained Approximation Error Lower Bound)}

        Let $A$ be a target matrix and $\hat{A}$ be an estimate. Let $S_A = \sr(A)$, $S_{\hat{A}} = \sr(\hat{A})$ with $S_{\hat{A}} \leqslant S_A$, and let $\sigma_1(\hat{A}) = \alpha \sigma_1(A)$ for some $\alpha \geq 0$, where $\sigma_1(\cdot)$ denotes the first (largest) singular value of the corresponding matrix. Then, the Frobenius distance is bounded below by:
        \[ 
        \|A - \hat{A}\|_F \geqslant \sigma_1(A)\left(\frac{\sqrt{S_A} - \sqrt{S_{\hat{A}}}}{1 + \sqrt{S_{\hat{A}}}}\right).
        \]
    \end{lemma}

    \begin{proof} $~~$\\
        By the definition of stable rank, $\sr(A) = \frac{\|A\|_F^2}{\sigma_1(A)^2}$, which implies $\|A\|_F = \sigma_1(A) \sqrt{\sr(A)}$. 
        By the reverse triangle inequality, we have:
        \[ \|A - \hat{A}\|_F \geqslant \left| \|A\|_F - \|\hat{A}\|_F \right|.\]
        Substituting the stable rank expressions yields:
        \[ \|A - \hat{A}\|_F \geqslant |\sigma_1(A)\sqrt{S_A} - \sigma_1(\hat{A})\sqrt{S_{\hat{A}}}| = \sigma_1(A)\left|\sqrt{S_A} - \alpha\sqrt{S_{\hat{A}}}\right|. \]
        We analyze this lower bound by considering two cases for $\alpha$: \\
        \textbf{Case 1 ($\alpha \leq 1$):} 
            \[ \|A - \hat{A}\|_F \geqslant \sigma_1(A)\left(\sqrt{S_A} - \sqrt{S_{\hat{A}}}\right).\]
        \textbf{Case 2 ($\alpha > 1$):} 
            \[ \|A - \hat{A}\|_F \geqslant \|A - \hat{A}\|_{op} \geqslant |\sigma_1(A) - \sigma_1(\hat{A})| = \sigma_1(A)(\alpha - 1).\]
        Combining both cases, we can construct a unified lower bound by finding the minimum possible error over all choices of $\alpha \geq 1$:
        \[\|A - \hat{A}\|_F \geqslant \sigma_1(A)\min_{\alpha\geqslant 1}\max \left(\sqrt{S_A} - \alpha\sqrt{S_{\hat{A}}}, \alpha - 1\right)\geqslant \sigma_1(A)\left(\frac{\sqrt{S_A} - \sqrt{S_{\hat{A}}}}{1 + \sqrt{S_{\hat{A}}}}\right).\]
        
        Comparing the bounds from both cases, we note that $\sqrt{S_A} - \sqrt{S_{\hat{A}}} \geqslant \frac{\sqrt{S_A} - \sqrt{S_{\hat{A}}}}{1 + \sqrt{S_{\hat{A}}}}$. Therefore, the global minimum over $\alpha$ yields: 
        \[ \|A - \hat{A}\|_F \geqslant \sigma_1(A)\left(\frac{\sqrt{S_A} - \sqrt{S_{\hat{A}}}}{1 + \sqrt{S_{\hat{A}}}}\right). \]
    \end{proof}


    This first lemma provides a fundamental lower bound on the approximation accuracy of a high-SR matrix by a low-SR one, establishing the intuition that this bound is primarily dependent on the difference between their stable ranks. This insight naturally leads to the conclusion that when dealing with <<difficult>> target data that intrinsically requires a high-SR adapter to be predicted accurately, our prediction accuracy will be fundamentally limited if we restrict ourselves to a low-SR adapter. \\


    \begin{lemma} \textbf{(Optimal SR-Constrained Approximation Error Bound)}

        Assuming the conditions of the previous lemma, let $A_R = \sum_{i=1}^R \sigma_i u_i v_i^T$ be the optimal unconstrained rank-$R$ approximation of $A$ with singular values $\sigma_1 \geqslant \dots \geqslant \sigma_R$, and let its stable rank be $\sr(A_R) = \frac{\sum_{i=1}^R \sigma_i^2}{\sigma_1^2}$. 
        For a given maximum stable rank $S \in [1, R]$, let $\mathcal{M}_{R, S}$ be the family of matrices $\hat{A}$ such that $\rank(\hat{A}) \leqslant R$ and $\sr(\hat{A}) \leqslant S$.
        Then the globally optimal approximation $\hat{A} \in \mathcal{M}_{R, S}$ minimizing $\|A - \hat{A}\|_F^2$ is given by $\hat{A} = \sum_{i=1}^R \hat{\sigma}_i u_i v_i^T$.
        If $\sr(A_R) \leqslant S$, then $\hat{\sigma}_i = \sigma_i$. If $\sr(A_R) > S$, the optimal singular values are:
        \[
        \hat{\sigma}_1 = \frac{\sigma_1 + \sqrt{S-1}\sqrt{\sum_{j=2}^R\sigma_j^2}}{S}, \quad \quad \hat{\sigma}_i = \gamma \sigma_i \text{ for } i \geqslant 2,
        \]
        where $\gamma = \frac{\sqrt{S-1}\hat{\sigma}_1}{\sqrt{\sum_{j=2}^R\sigma_j^2}}$. The minimal approximation error is exactly:
        \[
        \|A - \hat{A}\|_F^2 = \sum\limits_{i=R+1}^{D}\sigma_i^2 + \Delta_S,
        \]
        where the additional penalty $\Delta_S$ due to the stable rank constraint is:
        \[
        \Delta_S = \begin{cases} 0 & \text{if } ~ \sr(A_R) \leqslant S, \\ \frac{1}{S}\left(\sqrt{\sum\limits_{i=2}^R\sigma_i^2} - \sqrt{S - 1}\sigma_1\right)^2 & \text{if } ~ \sr(A_R) > S. \end{cases}
        \]
    \end{lemma}

    \begin{proof} $~~$\\
        By expanding the squared Frobenius distance, we obtain 
        $$\|A - \hat{A}\|_F^2 = \|A\|_F^2 + \|\hat{A}\|_F^2 - 2\text{Tr}(A^T \hat{A})$$
        The constraints on $\hat{A}$ (rank bounded by $R$ and stable rank bounded by $S$) depend entirely on its singular values $\hat{\sigma}_i$ and are invariant to its singular vectors. Since $\|A\|_F^2$ is constant and $\|\hat{A}\|_F^2 = \sum_{i=1}^R \hat{\sigma}_i^2$, minimizing the distance for any valid set of singular values requires maximizing the trace $\text{Tr}(A^T \hat{A})$. 
        By von Neumann's trace inequality,
        $$\text{Tr}(A^T \hat{A}) \leqslant \sum_{i=1}^R \sigma_i \hat{\sigma}_i$$
        and equality is strictly achieved if and only if the singular vectors of $\hat{A}$ perfectly align with the top $R$ singular vectors of $A$. Assigning any component of $\hat{A}$ to the tail vectors ($i > R$) of $A$ would pair it with a strictly smaller $\sigma_i$, decreasing the trace and increasing the error. Thus, we reduced the problem to finding values $\hat{\sigma}_1 \geqslant \dots \geqslant \hat{\sigma}_R \geqslant 0$ that minimize 
            $$f(\hat{\sigma}) = \sum_{i=1}^R (\sigma_i - \hat{\sigma}_i)^2$$ 
        subject to the stable rank constraint 
            $$g(\hat{\sigma}) = \sum_{i=2}^R \hat{\sigma}_i^2 - (S-1)\hat{\sigma}_1^2 \leqslant 0.$$

        This forms a convex optimization problem: the objective function $f(\hat{\sigma})$ is strictly convex (a paraboloid), and the feasible region defined by $g(\hat{\sigma}) \leqslant 0$ for $\hat{\sigma}_1 > 0$ forms a convex set. Slater's condition is naturally satisfied since there exists an interior point strictly satisfying the inequality e.g., $\hat{\sigma}_1 = 1$, $\hat{\sigma}_{i \geqslant 2} = 0$ yields $1-S < 0$ for $S>1$. Consequently, strong duality holds, and the Karush-Kuhn-Tucker (KKT) conditions are both necessary and sufficient for the global optimum.

        If $\sr(A_R) \leqslant S$, the unconstrained minimum $\hat{\sigma}_i = \sigma_i$ lies inside the feasible region, the Lagrange multiplier $\lambda$ is zero, and $\Delta_S = 0$.
        If $\sr(A_R) > S$, the unconstrained minimum lies outside the feasible region. By the complementary slackness condition ($\lambda \cdot g(\hat{\sigma}) = 0$), the optimal solution must lie strictly on the boundary. Thus, the constraint becomes active: $\sum_{i=2}^R \hat{\sigma}_i^2 = (S - 1)\hat{\sigma}_1^2$, with $\lambda > 0$.

        We define the Lagrangian for the active constraint over the top $R$ components:
        \[ \mathcal{L} = \frac{1}{2}(\hat{\sigma}_1 - \sigma_1)^2 + \frac{1}{2}\sum_{i=2}^R (\hat{\sigma}_i - \sigma_i)^2 + \frac{\lambda}{2} \left( \sum_{i=2}^R \hat{\sigma}_i^2 - (S-1)\hat{\sigma}_1^2 \right). \]
        Setting the partial derivatives with respect to the tail components ($i \geqslant 2$) to zero yields 
        $$\hat{\sigma}_i - \sigma_i + \lambda \hat{\sigma}_i = 0 \implies \hat{\sigma}_i = \gamma \sigma_i,$$ 
        where 
        $$\gamma = \frac{1}{1+\lambda} > 0.$$

        Instead of solving a non-linear system for $\lambda$ directly, we express the scaling factor $\gamma$ via $\hat{\sigma}_1$. Substituting $\hat{\sigma}_i = \gamma \sigma_i$ into the active constraint gives 
        $$\gamma^2 \sum_{i=2}^R \sigma_i^2 = (S-1)\hat{\sigma}_1^2.$$ 
        Letting $\|v\| = \sqrt{\sum_{i=2}^R \sigma_i^2}$, we obtain 
        $$\gamma = \frac{\sqrt{S-1}\hat{\sigma}_1}{\|v\|}.$$


        Substituting this exact dependency back into the original objective reduces the problem to a 1D unconstrained minimization over $\hat{\sigma}_1$:
        \[ f(\hat{\sigma}_1) = (\hat{\sigma}_1 - \sigma_1)^2 + \sum_{i=2}^R (\gamma \sigma_i - \sigma_i)^2 = (\hat{\sigma}_1 - \sigma_1)^2 + (\|v\| - \sqrt{S-1}\hat{\sigma}_1)^2. \]
        Setting $f'(\hat{\sigma}_1) = 0$ yields:
        \[ \hat{\sigma}_1 - \sigma_1 - \sqrt{S-1}(\|v\| - \sqrt{S-1}\hat{\sigma}_1) = 0 \implies \hat{\sigma}_1 = \frac{\sigma_1 + \sqrt{S-1}\|v\|}{S}. \]
        Since $\sr(A_R) > S$, we have $\|v\|^2 > (S-1)\sigma_1^2$, which strictly ensures $\hat{\sigma}_1 > \sigma_1$. This uncovers a fundamental property: to optimally enforce a lowered stable rank, one cannot merely shrink the tail components; the principal component must simultaneously be inflated. Furthermore, this guarantees that singular values remain properly ordered, as 
        $$\hat{\sigma}_2 = \gamma \sigma_2 = \frac{\sqrt{S-1}\sigma_2}{\|v\|}\hat{\sigma}_1 < \hat{\sigma}_1.$$

        Finally, plugging the optimal $\hat{\sigma}_1$ back into the objective gives the exact minimal squared error $\Delta_S$ introduced on the top $R$ components:
        \[ \Delta_S = \left(\frac{\sqrt{S-1}\|v\| - (S-1)\sigma_1}{S}\right)^2 + \left(\frac{\|v\| - \sqrt{S-1}\sigma_1}{S}\right)^2 = \frac{1}{S}(\|v\| - \sqrt{S-1}\sigma_1)^2. \]
        Adding the irreducible tail error $\sum_{i=R+1}^D \sigma_i^2$ provides the absolute minimal distance over the entire $\mathcal{M}_{R, S}$ family.
    \end{proof}

    Consequently, the second lemma does not merely establish an upper bound, but actively constructs the theoretically optimal approximation matrix. It provides the exact analytical formulation for the minimal error introduced when a low stable rank constraint is enforced along with a regular rank
       

  
    %%%%%%%%%%%%%%%%%%%%%%%%%%%%%%%%%%%%%%%%%%%%%%%%%%%%%%%%%%%%%%%%%%%%%%%%%%%%%%%%
    \subsection{Fully Connected Network Approximation}

    Having established the core matrix approximation mechanics, we can now evaluate the expressive power of LoRA within a multi-layer deep network. Zeng and Lee \cite{zeng2024expressivepowerlowrankadaptation} previously proved that rank-$r$ LoRA adapters can represent any rank-$r$ perturbation of pretrained weights across layers, establishing an existence proof for algebraic expressivity. However, their bound implies perfect capacity up to the nominal rank. 
    
    By extending their, we show that following the same compositional logic but imposing the empirically observed stable rank limits yields a significantly more restricted bound. Theorem 1 demonstrates that the approximation error at each layer is explicitly dictated by the stable rank penalty $\Delta_S$ derived in the lemma 2. Furthermore, because the network is deep, these layer-wise constraints do not simply sum; they are multiplicatively amplified by the spectral norms of subsequent layers, meaning an inability to model high-SR features in early layers exponentially degrades the output expressivity.

    
    %%%%%%%%%%%%%%%%%%%%%%%%%%%%%%%%%%%%%%%%%%%%%%%%%%%%%%%%%%%%%%%%%%%%%%%%%%%%%%%%
    \subsubsection{Notation}
    
    Let $FNN_{L,D}(.;(W_l)_{l=1}^L, (b_l)_{l=1}^L)$ be a fully connected neural network with $L$ layers of width $D$ and ReLU activation functions. $W_l\in \mathbb{R}^{D\times D}$ are the weight matrices and $b_l\in \mathbb{R}^D$ are the bias vectors.\\ 
    In this work we use:
    $$\text{target FNN  }~~ \overline{f}:=FNN_{\overline{L},D}(.;(\overline{W}_l)_{l=1}^L, (\overline{b}_l)_{l=1}^L), $$ 
    $$\text{frozen FNN }~~ f_0:= FNN_{L,D}(.;(W_l)_{l=1}^L, (b_l)_{l=1}^L),$$
    $$\text{adapted FNN }~~ f:=FNN_{L,D}(.;(W_l + \Delta W_l)_{l=1}^L, (\hat{b}_l)_{l=1}^L)$$
    where $b_l, \overline{b}_l, \hat{b}_l \in \mathbb{R}^D$, $W_l, \overline{W}_l, \Delta W_l \in \mathbb{R}^{D\times D}$, and $\Delta W_l = A_l B_l$, where $A_l\in \mathbb{R}^{D\times R}$ and $B_l\in \mathbb{R}^{R\times D}$ and $R\leqslant D$ is the nominal LoRA rank. Additionally, we assume $\overline{L} = L$.  

    %%%%%%%%%%%%%%%%%%%%%%%%%%%%%%%%%%%%%%%%%%%%%%%%%%%%%%%%%%%%%%%%%%%%%%%%%%%%%%%%
    \subsubsection{Expressive power theorem}
    
    Building upon the approximation framework and applying the matrix reconstruction residual from Lemma 2, we derive the following expressivity bound: \\

    \begin{theorem} $~$\\
        Let $x\in \mathbb{R}^d$ be the random input with covariance matrix $\Sigma = \mathbb{E}[xx^T]$.
        Let $E_l = \overline{W}_l - W_l$ be the required weight discrepancy at layer $l$ to match the target. Let $\sigma_{l,i}$ be the singular values of $E_l$, and let $E_{l,R}$ be its optimal unconstrained rank-$R$ approximation from Eckart-Young-Mirsky theorem. 
        Assume $\hat{b}_l = \overline{b}_l$. 
        
        Then there exist adapters $\Delta W_l$ of rank $\leqslant R$ and $\forall l \in\overline{1, L}~~\sr(\Delta W_l) \leqslant S$ such that the expected approximation error is bounded by:

        $$ \mathbb{E} \| f(\mathbf{x}) - \overline{f}(\mathbf{x}) \|_2 \leqslant \beta \sum_{l=1}^L \left( \prod_{k=l+1}^L \| \overline{W}_k \|_2 \right) \mathcal{T}_l $$

        Where $\beta$ accounts for the magnitude of the parameters and the input:
        $$ \beta = \max_{i \in [L]} \left( \underbrace{ \left( \prod_{j=1}^i \|\overline{W}_j\| \right) \sqrt{\text{tr}(\Sigma)} }_{\text{Input Term}} + \underbrace{ \sum_{j=1}^i \left( \prod_{k=j+1}^i \|\overline{W}_k\| \right) \|\overline{b}_j\| }_{\text{Bias Term}} \right) $$ 
        and $\mathcal{T}_l$ is the localized error derived from \textbf{Lemma 2}:
        $$ \mathcal{T}_l = \sqrt{ \sum_{j=R+1}^D \sigma_{l,j}^2 + \Delta_{S,l} } $$
        with the stable rank penalty $\Delta_{S,l}$ defined as:
        $$ \Delta_{S,l} = \begin{cases} 0 & \text{if } ~ \sr(E_{l,R}) \leqslant S \\ \frac{1}{S}\left(\sqrt{\sum\limits_{j=2}^R\sigma_{l,j}^2} - \sqrt{S - 1}\sigma_{l,1}\right)^2 & \text{if } ~ \sr(E_{l,R}) > S \end{cases} $$

    \end{theorem}


    \begin{proof}$~$

        1.  Linearization:
            Let $h_l$ be the output of layer $l$ in the adapted model, and $\overline{h}_l$ in the target model.
            $$ h_l = \mathbf{ReLU}((W_l + \Delta W_l) h_{l-1} + \overline{b}_l) $$
            $$ \overline{h}_l = \mathbf{ReLU}(\overline{W}_l \overline{h}_{l-1} + \overline{b}_l) $$
            We define the error vector 
            $$e_l = h_l - \overline{h}_l = \mathbf{ReLU}((W_l + \Delta W_l) h_{l-1}+ \overline{b}_l) - \mathbf{ReLU}(\overline{W}_l \overline{h}_{l-1} + \overline{b}_l) $$
            As $\mathbf{ReLU}$ is 1-Lipschitz
            $$\|e_l\|_2 \leqslant \| ((W_l + \Delta W_l) h_{l-1}+ \overline{b}_l) - (\overline{W}_l \overline{h}_{l-1} + \overline{b}_l)\|_2 $$

            Add and subtract $\overline{W}_l h_{l-1}$:
            $$ \|e_l\|_2 \leqslant \|(W_l + \Delta W_l - \overline{W}_l) h_{l-1} + \overline{W}_l (h_{l-1} - \overline{h}_{l-1})\|_2 $$
            Let $M_l = W_l + \Delta W_l - \overline{W}_l$ be the approximation residual matrix for layer $l$.
            $$ \|e_l\|_2 \leqslant \|M_l h_{l-1} + \overline{W}_l e_{l-1}\|_2 $$
            

        2.  Unrolling the recursion:
            By induction from $l=1$ to $L$ (with $e_0 = 0$):
            $$ \|e_L\|_2 \leqslant \sum_{l=1}^L \left\| \prod_{k=l+1}^L \overline{W}_k \right\|_2 \| M_l \|_2 \| h_{l-1} \|_2 $$

        3.  Constructing the Adapter $\Delta W_l$:
            We need to minimize $\| M_l \|_F = \| (W_l + \Delta W_l) - \overline{W}_l \|_F$.
            This is equivalent to approximating $E_l = \overline{W}_l - W_l$ with $\Delta W_l$.
            Thus we can apply \textbf{Lemma 2} and this yields the Frobenius error:
            $$ \| M_l \|_F = \mathcal{T}_l = \sqrt{ \sum_{j=R+1}^D \sigma_{l,j}^2 + \Delta_{S,l} } $$
            Since $\|M_l\|_2 \leqslant \|M_l\|_F$, we use $\mathcal{T}_l$ as the bound for the matrix norm.

        4.  Handling the intermediate norms:
            We look at $\mathbb{E} \|e_L\|_2$.
            The term depends on $\mathbb{E} \|h_{l-1}\|_2$. $\|h_{l-1}\|_2 \leqslant (\prod_{j=1}^{l-1} \|\overline{W}_j\|_2) \|\mathbf{x}\|_2$.
            For the expectation, $\mathbb{E}\|\mathbf{x}\|_2 \leqslant \sqrt{\mathbb{E}\|\mathbf{x}\|^2} = \sqrt{\text{tr}(\Sigma)}$.

            We want to bound $\| \overline{h}_i \|_2$.

                $$ \| \overline{h}_1 \| \leqslant \| \overline{W}_1 \| \| \mathbf{x} \| + \| \overline{b}_1 \| $$
                $$ \| \overline{h}_2 \| \leqslant \| \overline{W}_2 \| \| \overline{h}_1 \| + \| \overline{b}_2 \| $$
                Substitute the bound for Layer 1:
                $$ \| \overline{h}_2 \| \leqslant \| \overline{W}_2 \| ( \| \overline{W}_1 \| \| \mathbf{x} \| + \| \overline{b}_1 \| ) + \| \overline{b}_2 \| $$
                $$ \| \overline{h}_2 \| \leqslant \underbrace{\| \overline{W}_2 \| \| \overline{W}_1 \| \| \mathbf{x} \|}_{\text{Input Contribution}} + \underbrace{\| \overline{W}_2 \| \| \overline{b}_1 \| + \| \overline{b}_2 \|}_{\text{Accumulated Bias Contribution}} $$
            
                By induction, the magnitude of the activation at layer $i$ is bounded by:
                $$ \| \overline{h}_i \| \leqslant \left( \prod_{j=1}^i \| \overline{W}_j \| \right) \| \mathbf{x} \| + \sum_{j=1}^i \left( \prod_{k=j+1}^i \| \overline{W}_k \| \right) \| \overline{b}_j \| $$
                Taking the expectation and the maximum over layers we get the $\beta$ defined in theorem conditions

        5.  Gather up the terms:
            $$ \mathbb{E} \| f(\mathbf{x}) - \overline{f}(\mathbf{x}) \|_2 = \mathbb{E}[\|e_L\|_2] \leqslant  \beta \sum_{l=1}^L \left( \prod_{k=l+1}^L \| \overline{W}_k \|_2 \right) \mathcal{T}_l $$

        \end{proof}

    While this theorem relies on sequential worst-case norm accumulations, rendering the numerical bound loose in practical highly-dimensional settings, it successfully isolates the exact mechanic by which stable rank limits expressiveness. It reveals that the actual capacity of an adapter is not dictated by whether its nominal rank is large enough to contain the target matrix, as previous bounds imply, but by whether the optimizer can sufficiently distribute the energy across its singular values to capture the target's effective dimension. When the optimization process is intrinsically biased toward low-SR matrices, $\mathcal{T}_l$ explodes for high-SR targets.


    %%%%%%%%%%%%%%%%%%%%%%%%%%%%%%%%%%%%%%%%%%%%%%%%%%%%%%%%%%%%%%%%%%%%%%%%%%%%%%%%
    \subsection{Transformer Network Approximation}

    Extending our analysis from FNNs to Transformers introduces several layers of additional complexity that require careful mathematical treatment. Unlike FNNs with fixed affine transformations, the Transformer's attention mechanism computes weights dynamically based on the input through the softmax of Query-Key scores. This means the effective network transformation depends non-linearly on the input distribution. Furthermore, each transformer block possesses dual error propagation paths via the residual connections surrounding both the Multi-Head Attention (MHA) and FFN sub-layers, causing errors at layer $l$ to accumulate from both sources. 

    The sensitivity of the softmax function itself also amplifies parameter errors, a dynamic we encapsulate through a score sensitivity constant $C_{score,l}^h$. By carefully decomposing the local distortion $\mathcal{D}_l$ into FFN parameter errors and attention head distortions, we can establish how stable rank limits on the individual $Q, K, V,$ and $O$ projection matrices affects the Transformer's total expressive capacity.


    %%%%%%%%%%%%%%%%%%%%%%%%%%%%%%%%%%%%%%%%%%%%%%%%%%%%%%%%%%%%%%%%%%%%%%%%%%%%%%%%
    \subsubsection{Notation}
    
    Let
    $$TN_{L,D}(\cdot;((W_{Ql}^h, W_{Kl}^h, W_{Vl}^h, W_{Ol}^h)_{h=1}^H, (W_{1l}, W_{2l}, b_{1l}, b_{2l}))_{l=1}^L, W_o)$$
    be a transformer network with $L$ transformer blocks and an output linear layer parameterized by $W_o$. For theoretical feasibility we assume no dropout or normalisation operations. Each block is a composition of an $H$-head self-attention mechanism \cite{vaswani2023attentionneed} with weights $(W_{Ql}^h, W_{Kl}^h, W_{Vl}^h, W_{Ol}^h)_{h=1}^H$ and a 2-layer FFN with weights $(W_{1l}, W_{2l})$ and biases $(b_{1l}, b_{2l})$. Additionally, we assume all $W\in\mathbb{R}^{D\times D}, ~b\in \mathbb{R}^D$.
    We use the same prefix notation as before:
    $$\text{target TN  }~~ \overline{f} := TN_{L,D}(\cdot;((\overline{W_{Ql}^h}, \overline{W_{Kl}^h}, \overline{W_{Vl}^h},\overline{W_{Ol}^h})_{h=1}^H,(\overline{W_{1l}}, \overline{W_{2l}}, \overline{b_{1l}}, \overline{b_{2l}}))_{l=1}^L,\overline{W_o}), $$ 
    $$\text{frozen TN }~~ f_0 := TN_{L,D}(\cdot;((W_{Ql}^h, W_{Kl}^h, W_{Vl}^h, W_{Ol}^h)_{h=1}^H, (W_{1l}, W_{2l}, b_{1l}, b_{2l}))_{l=1}^L, W_o),$$
    $$\text{adapted TN }~~ f := TN_{L,D}(\cdot;((W_{Ql}^h + \Delta W_{Ql}^h, W_{Kl}^h + \Delta W_{Kl}^h, W_{Vl}^h + \Delta W_{Vl}^h, W_{Ol}^h + \Delta W_{Ol}^h)_{h=1}^H, $$ 
    $$(W_{1l} + \Delta W_{1l}, W_{2l} + \Delta W_{2l}, \widehat{b_{1l}}, \widehat{b_{2l}}))_{l=1}^L, W_o + \Delta W_o)$$
    
    For a matrix $G\in \mathbb{R}^{D\times D}$, let its optimal rank-$R$ approximation be $G_R$ with singular values $\sigma_1(G) \geqslant \dots \geqslant \sigma_R(G)$. Let $\mathcal{T}(G)$ be the localized SR error from \textbf{Lemma 2}:
    $$\mathcal{T}(G) = \sqrt{\sum\limits_{i=R+1}^{D}\sigma_i^2(G) + \Delta_S(G)}$$
    where $\Delta_S(G) = 0$ if $\sr(G_R) \leqslant S$, and $\Delta_S(G) = \frac{1}{S}\left(\sqrt{\sum_{i=2}^R\sigma_i^2(G)} - \sqrt{S - 1}\sigma_1(G)\right)^2$ if $\sr(G_R) > S$.




    %%%%%%%%%%%%%%%%%%%%%%%%%%%%%%%%%%%%%%%%%%%%%%%%%%%%%%%%%%%%%%%%%%%%%%%%%%%%%%%%
    \subsubsection{Expressive power of transformer network}

    Following the block decomposition approach used in \cite{yun2019transformers} and extending the rigorous Lipschitz analysis from \cite{virmaux2018lipschitz}, we formulate the following bound tracking the accumulation of stable rank constraints in a self-attention architecture:

    \begin{theorem} $~$
    
        Let $x\in \mathbb{R}^{D\times N}$ be the random input sequence. Assume $\overline{b_{\cdot\cdot}} = \widehat{b_{\cdot\cdot}}$. Let $z_l\in \mathbb{R}^{D\times N}$ be the output of $l$ block in the adapted model and $\overline{z_l}$ be the output of the same block in the target model. Consider $z_0 = \overline{z_0} = x$. \\
        Let us define per matrix discrepancies as follows:
        $$\forall l\in\overline{1,L},~h\in\overline{1,H},~ \forall ~ role\in \{Q, K, V, O\}: ~~\mathcal{T}_{role, l}^h =\mathcal{T}(W_{role, l}^h) = \|W_{role, l}^h - \overline{W_{role, l}^h}\|_F$$
        $$\forall l\in\overline{1,L},~ \forall ~ role \in \{1, 2\}: ~~\mathcal{T}_{role, l} =\mathcal{T}(W_{role, l}) = \|W_{role, l}^h - \overline{W_{role, l}}\|_F$$
        $$\mathcal{T}_o =\mathcal{T}(W_o) = \|W_o - \overline{W_o}\|_F$$
        Additionally we define the following bound constants:
        $$\forall l\in\overline{1,L},~h\in\overline{1,H},~ \forall ~ role\in \{Q, K, V, O\}: ~~\mu_{role, l}^h = \|\overline{W_{role, l}^h}\|_2 + \mathcal{T}_{role, l}^h$$
        $$\forall l\in\overline{1,L},~ \forall ~ role \in \{1, 2\}: ~~\mu_{role, l} = \|\overline{W_{role, l}}\|_2 + \mathcal{T}_{role, l}$$
        $$\mu_o = \|\overline{W_o}\|_2 + \mathcal{T}_o$$
        $$\beta_0 = \sqrt{\mathbb{E}[\|x\|^2_F]}\text{, $~~$ for l $>$ 0: } \beta_l=\beta_{l-1}\cdot (1+\mu_{2,l}\mu_{1,l})\left(1 + \sum\limits_{h=1}^H\mu_{O,l}^h\mu_{V,l}^h \right)$$
        $$\beta = \beta_L$$

        Then there exists a $\leq S$ stable rank LoRA adapter configuration such that
        $$\mathbb{E}\|f(x) - \overline{f}(x)\|_F \leqslant \beta\sum\limits_{l=1}^{L}\left(\prod\limits_{k=l+1}^L \gamma_k \right)\mathcal{D}_l $$
        where:

        \begin{enumerate}
            \item $\gamma_k$ is the Lipschitz constant of target layer k with respect to the input:
                $$\gamma_k = \|\overline{W_{2,k}^h}\|\|\overline{W_{1,k}^h}\|\left(1 + \sum\limits_{h=1}^H\Psi_k^h\right)$$
                $$ \Psi_k^h = \left\|\overline{W_{O,k}^h}\right\|_2 \left( \left\|\overline{W_{V,k}^h}\right\|_2 + \frac{\lambda_{sm} \beta^2}{\sqrt{D}} \left\|\overline{W_{V,k}^h}\right\|_2 \left(\mu_{Q,k}^h \left\|\overline{W_{K,k}^h}\right\|_2 + \left\|\overline{W_{K,k}^h}\right\|_2 \left\|\overline{W_{Q,k}^h}\right\|_2\right) \right) $$

            \item $\mathcal{D}_l$ is the Local Distortion at layer $l$, defined as:
                    $$ \mathcal{D}_l = \mathcal{T}_{2,l} + \|\overline{W_{2,l}}\|_2 \mathcal{T}_{1,l} + \|\overline{W_{2,l}}\|_2 \|\overline{W_{1,l}}\|_2 \sum_{h=1}^H \mathcal{D}_{att, l}^h $$
                With the per-head attention distortion $\mathcal{D}_{att, l}^h$:
                $$ \mathcal{D}_{att, l}^h = \mu_{V,l}^h \mathcal{T}_{O,l}^h + \|\overline{W_{O,l}^h}\|_2 \mathcal{T}_{V,l}^h + C_{score,l}^h (\mu_{Q,l}^h \mathcal{T}_{K,l}^h + \|\overline{W_{K,l}^h}\|_2 \mathcal{T}_{Q,l}^h) $$
                And the score sensitivity constant:
                    $$ C_{score,l}^h = \frac{\lambda_{sm}}{\sqrt{D}} \|\overline{W_{O,l}^h}\|_2 \|\overline{W_{V,l}^h}\|_2 \beta $$
        \end{enumerate}               

    \end{theorem}

    \begin{proof}$~$

        Let $E_l = z_l - \overline{z}_l$ be the error at the output of layer $l$.

        The output of layer $l$ is $z_l = \text{FFN}_l(z_{l-1} + \text{MHA}_l(z_{l-1}))$. Let $Y_l = z_{l-1} + \text{MHA}_l(z_{l-1})$ be the intermediate state.
        $$ E_l = \text{FFN}_l(Y_l) - \overline{\text{FFN}_l}(\overline{Y_l}) $$
        $$ \|E_l\|_F \le \underbrace{\left\| \text{FFN}_l(Y_l) - \overline{\text{FFN}_l}(Y_l) \right\|_F}_{\textbf{FFN Parameter Error}} + \underbrace{\left\| \overline{\text{FFN}_l}(Y_l) - \overline{\text{FFN}_l}(\overline{Y_l}) \right\|_F}_{\textbf{FFN Input Error}} $$

        The FFN part of TN layer $l$ is $W_{2,l} \text{ReLU}(W_{1,l} x)$. Then the \textbf{FFN parameter error}:
        $$\left\|W_{2,l} \text{ReLU}(W_{1,l} Y_l) - \overline{W_{2,l}} \text{ReLU}(\overline{W_{1,l}} Y_l)\right\|_F \leqslant \left\|(W_{2,l} - \overline{W_{2,l}})\text{ReLU}(W_{1,l} Y_l)\right\|_F + \left\|\overline{W_{2,l}}\right\|_2 \left\|\text{ReLU}(W_{1,l} Y_l) - \text{ReLU}(\overline{W_{1,l}} Y_l)\right\|_F \leqslant $$
        $$ \leqslant \mathcal{T}_{2,l} \|Y_l\|_F + \left\|\overline{W_{2,l}}\right\|_2 \cdot \left\|(W_{1,l} - \overline{W_{1,l}}) Y_l\right\|_F \leqslant \beta \mathcal{T}_{2,l} + \left\|\overline{W_{2,l}}\right\|_2 \beta \mathcal{T}_{1,l} $$

        The \textbf{FFN Input Error} is bounded by the following
        $$\left\|\overline{W_{2,l}}(\text{ReLU}(\overline{W_{1,l}} Y_l) - \text{ReLU}(\overline{W_{1,l}} \overline{Y_l}))\right\|_F \leqslant \left\|\overline{W_{2,l}}\right\|_2 \cdot \left\|\overline{W_{1,l}}\right\|_2 \cdot \left\|Y_l - \overline{Y_l}\right\|_F $$

        Therefore the combined total FFN error bound:
        $$ \|E_l\|_F \leqslant \beta (\mathcal{T}_{2,l} + \left\|\overline{W_{2,l}}\right\|_2 \mathcal{T}_{1,l}) + \left\|\overline{W_{2,l}}\right\|_2 \cdot \left\|\overline{W_{1,l}}\right\|_2 \cdot \left\|Y_l - \overline{Y_l}\right\|_F $$

        Now we analyze the Attention error $\|Y_l - \overline{Y_l}\|_F$. Since $Y_l = z_{l-1} + \sum_h \text{Head}_h^l$, and $E_{l-1} = z_{l-1} - \overline{z_{l-1}}$:
        $$ \left\|Y_l - \overline{Y_l}\right\|_F \leqslant \|E_{l-1}\|_F + \sum_{h=1}^H \left\| \text{Head}_h^l(z_{l-1}) - \overline{\text{Head}_h^l}(\overline{z_{l-1}}) \right\|_F $$

        A single head output: 
        $$\text{Head}_h^l(z_{l-1}) = W_{O,l}^h V_l^h (A_l^h)^T$$ 
        where 
        $$V_l^h = W_{V, l}^h z_{l-1}, ~~ A_l^h = \text{Softmax}(S_l^h), ~~S_l^h = \frac{1}{\sqrt{D}} (K_l^h)^T Q_l^h, ~~K_l^h=W_{K,l}^h z_{l-1},~~ Q_l^h=W_{Q,l}^h z_{l-1}$$

        As with the FFN part, we decompose the error $\|\text{Head} - \overline{\text{Head}}\|_F$ into \textbf{Head parameter Error} (functions of $\mathcal{T}$) and \textbf{Head Input Error} (functions of $E_{l-1}$).

        Output projection $W_{O,l}^h$:
        $$ W_{O,l}^h V_l^h (A_l^h)^T - \overline{W_{O,l}^h} \overline{V_l^h} (\overline{A_l^h})^T = \left(W_{O,l}^h - \overline{W_{O,l}^h}\right) V_l^h (A_l^h)^T + \overline{W_{O,l}^h} \left(V_l^h (A_l^h)^T - \overline{V_l^h} (\overline{A_l^h}^T)\right) $$

        For the first term the following holds:
        $$\left\|\left(W_{O,l}^h - \overline{W_{O,l}^h}\right) V_l^h (A_l^h)^T \right\|_F \leqslant \left\|W_{O,l}^h - \overline{W_{O,l}^h}\right\|_F\cdot \left\|V_l^h (A_l^h)^T\right\|_2 \leqslant \left\|W_{O,l}^h - \overline{W_{O,l}^h}\right\|_F\cdot \left\|V_l^h (A_l^h)^T\right\|_F \leqslant $$
        $$\leqslant \mathcal{T}_{O,l}^h \cdot \left\|V_l^h\right\|_F \cdot \left\|(A_l^h)^T\right\|_2\leqslant \mathcal{T}_{O,l}^h \cdot \left\|W_{V,l}^h z_{l-1}\right\|_F \cdot 1 \leqslant \mathcal{T}_{O,l}^h \cdot \left\|W_{V,l}^h\right\|_2\cdot \|z_{l-1}\|_F \leqslant \mathcal{T}_{O,l}^h \cdot \mu_{V,l}^h\cdot \beta$$
        
        For the second term:
        $$\left\|\overline{W_{O,l}^h} \left(V_l^h (A_l^h)^T - \overline{V_l^h} (\overline{A_l^h})^T\right)\right\| \leqslant \left\|\overline{W_{O,l}^h}\right\|_2 \left\|V_l^h (A_l^h)^T - \overline{V_l^h} (\overline{A_l^h})^T\right\|_F$$
        $$ V_l^h (A_l^h)^T - \overline{V_l^h} (\overline{A_l^h})^T = (V_l^h - \overline{V_l^h}) (A_l^h)^T + \overline{V_l^h} (A_l^h - (\overline{A_l^h})^T) $$
        $V$ matrices difference:
        $$V_l^h - \overline{V_l^h} = (W_{V,l}^h - \overline{W_{V,l}^h})z_{l-1} + \overline{W_{V,l}^h}(z_{l-1} - \overline{z_{l-1}})$$
        $$\left\|V_l^h - \overline{V_l^h}\right\|_F\leqslant \left\|(W_{V,l}^h - \overline{W_{V,l}^h})z_{l-1}\right\|_F + \left\|\overline{W_{V,l}^h}(z_{l-1} - \overline{z_{l-1}})\right\|_F\leqslant \mathcal{T}_{V,l}^h \beta + \left\|\overline{W_{V,l}^h}\right\|_2\cdot\|E_{l-1}\|_F$$
        $A$ matrices difference:

        $$\left\|A_l^h - \overline{A_l^h}\right\|_F \leqslant \lambda_{sm} \left\|S_l^h - \overline{S_l^h}\right\|_F $$

        Additionally $\left\|\overline{V_l^h}\right\|_F \leqslant \left\|\overline{W_{V,l}^h}\right\|_2 \cdot \beta ~~$, thus:

        $$\left\|\overline{V_l^h} (A_l^h - (\overline{A_l^h})^T)\right\|\leqslant \lambda_{sm} \cdot \beta \cdot\left\|\overline{W_{V,l}^h}\right\|_2 \cdot \left\|S_l^h - \overline{S_l^h}\right\|_F$$

        Now we need to bound the score matrices $S$ difference:
        $$ S_l^h - \overline{S_l^h} = \frac{1}{\sqrt{D}} ((K_l^h)^T Q_l^h - (\overline{K_l^h})^T \overline{Q_l^h}) = \frac{1}{\sqrt{D}}\left((K_l^h - \overline{K_l^h})^T Q_l^h + (\overline{K_l^h})^T(Q_l^h - \overline{Q_l^h})\right)$$

        $Q$ matrices difference:
         $$Q_l^h - \overline{Q_l^h} = (W_{Q, l}^h - \overline{W_{Q, l}^h}) z_{l-1} + \overline{W_{Q, l}^h} (z_{l-1} - \overline{z_{l-1}})$$
         $$\left\|Q_l^h - \overline{Q_l^h}\right\|_F \leqslant \mathcal{T}_{Q,l}^h \cdot \beta + \left\|\overline{W_{Q,l}^h}\right\|_2 \cdot \left\|E_{l-1}\right\|_F$$

        $K$ matrices difference:
         $$K_l^h - \overline{K_l^h} = (W_{K,l}^h - \overline{W_{K,l}^h}) z_{l-1} + \overline{W_{K,l}^h} (z_{l-1} - \overline{z_{l-1}})$$
         $$\left\|K_l^h - \overline{K_l^h}\right\|_F \leqslant \mathcal{T}_{K,l}^h \cdot \beta + \left\|\overline{W_{K,l}^h}\right\|_2 \cdot \left\|E_{l-1}\right\|_F$$

        Now we substitute back step by step:
        $$\left\|S_l^h - \overline{S_l^h}\right\|_F \leqslant \frac{1}{\sqrt{D}}\left(\left\|(K_l^h - \overline{K_l^h})^T Q_l^h\right\|_F + \left\|(\overline{K_l^h})^T(Q_l^h - \overline{Q_l^h})\right\|_F\right) \leqslant$$
        $$\leqslant \frac{1}{\sqrt{D}}\left(\left\|K_l^h - \overline{K_l^h}\right\|_F \left\|W_{Q,l}^h\right\|_2 \left\|z_{l-1}\right\|_F +  \left\|\overline{W_{K,l}^h}\right\|_2 \left\|\overline{z_{l-1}}\right\|_F \left\|Q_l^h - \overline{Q_l^h}\right\|_F\right) \leqslant $$
        $$\leqslant \frac{\beta}{\sqrt{D}}\left(\left\|K_l^h - \overline{K_l^h}\right\|_F \mu_{Q,l}^h +  \left\|\overline{W_{K,l}^h}\right\|_2  \left\|Q_l^h - \overline{Q_l^h}\right\|_F\right) \leqslant$$
        $$ \leqslant \frac{\beta}{\sqrt{D}}\left(\mu_{Q,l}^h \mathcal{T}_{K,l}^h \beta + \left\|\overline{W_{K,l}^h}\right\|_2 \mathcal{T}_{Q,l}^h \beta \right) + \frac{\beta}{\sqrt{D}}\left(\mu_{Q,l}^h \left\|\overline{W_{K,l}^h}\right\|_2 + \left\|\overline{W_{K,l}^h}\right\|_2 \left\|\overline{W_{Q,l}^h}\right\|_2\right) \left\|E_{l-1}\right\|_F$$
        $$\left\|\overline{V_l^h} (A_l^h - (\overline{A_l^h})^T)\right\|_F \leqslant $$
        $$\leqslant \lambda_{sm} \beta \left\|\overline{W_{V,l}^h}\right\|_2 \left[ \frac{\beta}{\sqrt{D}}\left(\mu_{Q,l}^h \mathcal{T}_{K,l}^h \beta + \left\|\overline{W_{K,l}^h}\right\|_2 \mathcal{T}_{Q,l}^h \beta \right) + \frac{\beta}{\sqrt{D}}\left(\mu_{Q,l}^h \left\|\overline{W_{K,l}^h}\right\|_2 + \left\|\overline{W_{K,l}^h}\right\|_2 \left\|\overline{W_{Q,l}^h}\right\|_2\right) \left\|E_{l-1}\right\|_F \right]$$
        $$\left\|V_l^h (A_l^h)^T - \overline{V_l^h} (\overline{A_l^h})^T\right\|_F \leqslant \left\|V_l^h - \overline{V_l^h}\right\|_F \left\|(A_l^h)^T\right\|_2 + \left\|\overline{V_l^h} (A_l^h - (\overline{A_l^h})^T)\right\|_F\leqslant $$
        $$\leqslant  \mathcal{T}_{V,l}^h \beta + \left\|\overline{W_{V,l}^h}\right\|_2 \left\|E_{l-1}\right\|_F + \frac{\lambda_{sm} \beta^2}{\sqrt{D}} \left\|\overline{W_{V,l}^h}\right\|_2 \left(\mu_{Q,l}^h \mathcal{T}_{K,l}^h \beta + \left\|\overline{W_{K,l}^h}\right\|_2 \mathcal{T}_{Q,l}^h \beta \right) + $$
        $$ + \frac{\lambda_{sm} \beta^2}{\sqrt{D}} \left\|\overline{W_{V,l}^h}\right\|_2 \left(\mu_{Q,l}^h \left\|\overline{W_{K,l}^h}\right\|_2 + \left\|\overline{W_{K,l}^h}\right\|_2 \left\|\overline{W_{Q,l}^h}\right\|_2\right) \left\|E_{l-1}\right\|_F$$
        Getting back to the head error:
        $$\left\|\text{Head}_h^l(z_{l-1}) - \overline{\text{Head}_h^l}(\overline{z_{l-1}})\right\|_F \leqslant \mathcal{T}_{O,l}^h \mu_{V,l}^h \beta + \left\|\overline{W_{O,l}^h}\right\|_2 \left\|V_l^h (A_l^h)^T - \overline{V_l^h} (\overline{A_l^h})^T\right\|_F\leqslant$$
        $$\leqslant\mathcal{T}_{O,l}^h \mu_{V,l}^h \beta + \left\|\overline{W_{O,l}^h}\right\|_2 \left[ \mathcal{T}_{V,l}^h \beta + \frac{\lambda_{sm} \beta^2}{\sqrt{D}} \left\|\overline{W_{V,l}^h}\right\|_2 \left(\mu_{Q,l}^h \mathcal{T}_{K,l}^h \beta + \left\|\overline{W_{K,l}^h}\right\|_2 \mathcal{T}_{Q,l}^h \beta \right) \right] + $$
        $$ + \left\|\overline{W_{O,l}^h}\right\|_2 \left[ \left\|\overline{W_{V,l}^h}\right\|_2 + \frac{\lambda_{sm} \beta^2}{\sqrt{D}} \left\|\overline{W_{V,l}^h}\right\|_2 \left(\mu_{Q,l}^h \left\|\overline{W_{K,l}^h}\right\|_2 + \left\|\overline{W_{K,l}^h}\right\|_2 \left\|\overline{W_{Q,l}^h}\right\|_2\right) \right] \left\|E_{l-1}\right\|_F$$
        
        Grouping the error into \textbf{Head parameter Error} ($\mathcal{D}_{att, l}^h$) and \textbf{Head Input Error} ($\Psi_l^h$):
        $$\left\|\text{Head}_h^l(z_{l-1}) - \overline{\text{Head}_h^l}(\overline{z_{l-1}})\right\|_F \leqslant \beta \mathcal{D}_{att, l}^h + \Psi_l^h \left\|E_{l-1}\right\|_F$$

        where
        $$\mathcal{D}_{att, l}^h = \mathcal{T}_{O,l}^h \mu_{V,l}^h + \left\|\overline{W_{O,l}^h}\right\|_2 \mathcal{T}_{V,l}^h + \frac{\lambda_{sm} \beta}{\sqrt{D}} \left\|\overline{W_{O,l}^h}\right\|_2 \left\|\overline{W_{V,l}^h}\right\|_2 \left(\mu_{Q,l}^h \mathcal{T}_{K,l}^h \beta + \left\|\overline{W_{K,l}^h}\right\|_2 \mathcal{T}_{Q,l}^h \beta \right)$$
        
        $$\Psi_l^h = \left\|\overline{W_{O,l}^h}\right\|_2 \left( \left\|\overline{W_{V,l}^h}\right\|_2 + \frac{\lambda_{sm} \beta^2}{\sqrt{D}} \left\|\overline{W_{V,l}^h}\right\|_2 \left(\mu_{Q,l}^h \left\|\overline{W_{K,l}^h}\right\|_2 + \left\|\overline{W_{K,l}^h}\right\|_2 \left\|\overline{W_{Q,l}^h}\right\|_2\right) \right)$$

        The output of the multi-head attention block combined with the residual connection is denoted as $Y_l$:
        $$Y_l = z_{l-1} + \sum_{h=1}^H \text{Head}_h^l(z_{l-1})$$

        The difference before the FFN layer is bounded by summing the individual head errors:
        $$\|Y_l - \overline{Y_l}\|_F = \left\|z_{l-1} - \overline{z_{l-1}} + \sum_{h=1}^H \left( \text{Head}_h^l(z_{l-1}) - \overline{\text{Head}_h^l}(\overline{z_{l-1}}) \right)\right\|_F$$
        $$\|Y_l - \overline{Y_l}\|_F \leqslant \|E_{l-1}\|_F + \sum_{h=1}^H \left( \beta \mathcal{D}_{att, l}^h + \Psi_l^h \|E_{l-1}\|_F \right)$$
        $$\|Y_l - \overline{Y_l}\|_F \leqslant \beta \sum_{h=1}^H \mathcal{D}_{att, l}^h + \left( 1 + \sum_{h=1}^H \Psi_l^h \right) \|E_{l-1}\|_F$$
        
        As mentioned above, the output error of the FFN is bounded by
        $$ \|E_l\|_F \leqslant \beta (\mathcal{T}_{2,l} + \left\|\overline{W_{2,l}}\right\|_2 \mathcal{T}_{1,l}) + \left\|\overline{W_{2,l}}\right\|_2 \cdot \left\|\overline{W_{1,l}}\right\|_2 \cdot \left\|Y_l - \overline{Y_l}\right\|_F $$

        Substitute the bound for $\|Y_l - \overline{Y_l}\|_F$ into the equation:
        $$\|E_l\|_F \leqslant \beta \mathcal{T}_{2,l} + \beta \left\|\overline{W_{2,l}}\right\|_2 \mathcal{T}_{1,l} + \left\|\overline{W_{2,l}}\right\|_2 \left\|\overline{W_{1,l}}\right\|_2 \left[ \beta \sum_{h=1}^H \mathcal{D}_{att, l}^h + \left( 1 + \sum_{h=1}^H \Psi_l^h \right) \|E_{l-1}\|_F \right]$$
        
        Group the terms into the local layer distortion $\mathcal{D}_l$ and the layer Lipschitz amplification $\Lambda_l$:
        $$\|E_l\|_F \leqslant \beta \underbrace{\left[ \mathcal{T}_{2,l} + \left\|\overline{W_{2,l}}\right\|_2 \mathcal{T}_{1,l} + \left\|\overline{W_{2,l}}\right\|_2 \left\|\overline{W_{1,l}}\right\|_2 \sum_{h=1}^H \mathcal{D}_{att, l}^h \right]}_{\mathcal{D}_l} + \underbrace{\left[ \left\|\overline{W_{2,l}}\right\|_2 \left\|\overline{W_{1,l}}\right\|_2 \left( 1 + \sum_{h=1}^H \Psi_l^h \right) \right]}_{\Lambda_l} \|E_{l-1}\|_F$$
        $$\|E_l\|_F \leqslant \beta \mathcal{D}_l + \Lambda_l \|E_{l-1}\|_F$$
        By unrolling the recurrence, considering $\|E_0\|_F = 0$:

        $$\mathbb{E} \|E_L\|_F \leqslant \mathbb{E} \left[ \sum_{l=1}^L \left( \prod_{k=l+1}^L \Lambda_k \right) \beta \mathcal{D}_l \right]$$
        $$\mathbb{E} \|f(X) - \tilde{f}(X)\|_F \leqslant \beta \sum_{l=1}^L \left( \prod_{k=l+1}^L \Lambda_k \right) \mathcal{D}_l$$

    \end{proof}


    While the explicit algebraic expansion relies on sequential worst-case multiplicative error accumulations, rendering the numerical bound loose in practical high-dimensional settings, it establishes a rigorous theoretical linkage between stable rank and deep network capacity. We have successfully traced the SR-constrained matrix error $\mathcal{T}$ from a single linear projection through the highly non-linear, dual-path architecture of a multi-head transformer. 
    
    The theorem formalizes how applying LoRA to the Query, Key, Value, and Output projection matrices restricts the adaptability of the attention head according to the stable rank limit. Even though the frozen base matrices retain their full rank, steering the attention mechanism toward a significantly different target function might require a high-SR matrix. If the downstream data intrinsically demands high stable rank updates to form new associations between tokens, an optimization trajectory that implicitly biases toward low sr solutions will fundamentally cap the Transformer's representational accuracy. This mathematical bottleneck suggests a strong hypothesis for the performance plateaus often observed empirically at high nominal ranks: real expressivity scales primarily with stable rank rather than algebraic rank, meaning that simply increasing the parameter budget $r$ yields diminishing returns if the effective dimension of the update remains constricted by the optimization algorithm.



    %%%%%%%%%%%%%%%%%%%%%%%%%%%%%%%%%%%%%%%%%%%%%%%%%%%%%%%%%%%%%%%%%%%%%%%%%%%%%%%%
    \subsection{Generalization Analysis}

    Understanding generalization -- why a model performs well on unseen data despite being most of the times massively overparameterized -- is a central pursuit in deep learning theory. The work of Kratsios et al. \cite{kratsios2025sharpgeneralizationboundsfoundation} recently established generalization bounds for low-rank adapters in terms of rank and sample complexity. Their upper bound scales as $\mathcal{O}(\sqrt{r/N})$, where $r$ is the nominal adapter rank and $N$ is the sample size. This bound, derived from worst-case estimates of rank-$r$ matrices, theoretically suggests that generalization degrades proportionally as one increases the nominal rank. 
    
    However, this analysis relies purely on structural constraints and does not account for the continuous spectral structure or energy distribution of the learned adapters. The gap between Kratsios et al.'s bound and empirical performance (where larger ranks often do not overfit as aggressively as predicted) stems directly from their worst-case analysis assumptions. Bounds relying purely on algebraic rank are mathematically tight only when all $r$ singular values are roughly equal ($\sigma_1 \approx \sigma_2 \approx \cdots \approx \sigma_r$), representing the maximum entropy distribution under a rank constraint. As LoRA optimization empirically often collapses into highly skewed distributions with low stable rank, evaluating complexity purely by nominal rank vastly overestimates the capacity of the model and the resulting generaization gap is pessimistically high.
    
    In this section, we rigorously prove that the generalization bound can be tightened by indtriducing a stable rank rank constraint. Specifically, we prove that our proxy function for generalization gap scales with $\sqrt{\sr}$ rather than $\sqrt{\rank}$ leveraging the similar PAC-Bayesian pipeline.



    %%%%%%%%%%%%%%%%%%%%%%%%%%%%%%%%%%%%%%%%%%%%%%%%%%%%%%%%%%%%%%%%%%%%%%%%%%%%%%%%
    \subsubsection{PAC-Bayesian Framework Integration}

    The PAC-Bayesian framework \cite{mcallester1999pac, mcallester2003simplified, shawe1997pac} provides an elegant and powerful tool for deriving generalization bounds. Unlike VC-dimension or Rademacher complexity bounds---which evaluate the combinatorial capacity of an entire hypothesis class---PAC-Bayes provides data-dependent bounds that can natively incorporate prior knowledge about the hypothesis space. This is uniquely crucial for analyzing fine-tuning.

    The key mechanism bounds the generalization gap through the Kullback-Leibler (KL) divergence $KL(\mathcal{Q} || \mathcal{P})$ between a reference distribution $\mathcal{P}$ and a learned distribution $\mathcal{Q}$ over the network's weights. In the PAC-Bayes literature, $\mathcal{P}$ and $\mathcal{Q}$ are conventionally termed the <<prior>> and <<posterior>> distributions. However, it is important to note that these do not represent a strict Bayesian inference update (i.e., $P(\theta|D) \propto P(D|\theta)P(\theta)$). Rather, $\mathcal{P}$ can be any distribution chosen \textit{before} observing the training data, and $\mathcal{Q}$ can be any distribution chosen afterward. 
    
    For evaluating LoRA, we define our prior $\mathcal{P}$ as an isotropic Gaussian centered on the frozen pretrained weight matrices $W_{0,t}$ at each layer $t$: $\mathcal{P} = \prod_t \mathcal{N}(W_{0,t}, \tau_t^2 I)$, where $\tau_t$ is a layer-wise variance hyperparameter and $I$ is the identity matrix. For the posterior $\mathcal{Q}$, we center the distribution on the final adapted weights $W_{0,t} + \Delta W_t$, where $\Delta W_t$ is the learned LoRA adapter, yielding $\mathcal{Q} = \prod_t \mathcal{N}(W_{0,t} + \Delta W_t, \tau_t^2 I)$. Under this formulation, the KL divergence cleanly reduces to a sum of squared Frobenius norms: $KL(\mathcal{Q} || \mathcal{P}) = \sum_t \frac{\|\Delta W_t\|_F^2}{2\tau_t^2}$.

    While Neyshabur et al. \cite{neyshabur2017pac} originally utilized this PAC-Bayesian formulation to analyze standard, full-rank feed-forward networks, adapting their framework to LoRA fine-tuning naturally isolates the Frobenius norm of the adapters themselves: $\|\Delta W_t\|_F^2$. If one were to naively bound the capacity of these adapters, the standard approach would rely on their algebraic rank $r$, substituting the Frobenius norm with the worst-case volumetric bound $\|\Delta W_t\|_F^2 \leq \Lambda_t^2 r$ (where $\Lambda_t$ bounds the spectral norm). However, as Frobenius norm, by definition of stable rank, decomposes as $\|\Delta W_t\|_F^2 = \|\Delta W_t\|_{op}^2 \cdot \sr(\Delta W_t)$, we can directly replace this by $\|\Delta W_t\|_F^2 \leq \Lambda_t^2 S_l$, where $S_l$ is the stable rank constraint for layer $l$. 

    

    %%%%%%%%%%%%%%%%%%%%%%%%%%%%%%%%%%%%%%%%%%%%%%%%%%%%%%%%%%%%%%%%%%%%%%%%%%%%%%%%
    \subsubsection{Notation}
    
    In this section, we use a setting similar to the one described in \cite{kratsios2025sharpgeneralizationboundsfoundation} and expressivity section above. Consider a depth-$T$ feed-forward neural network $f_{W_0}$ with 1-Lipschitz activation functions $\sigma(\cdot)$ such that $\sigma(0) = 0$. The pre-trained network is parameterized by weight matrices $W_{0,t} \in \mathbb{R}^{d_{t-1}\times d_{t}}$ for each layer $t$.
    LoRA is applied via matrices $\Delta W_t = B_t(A_t)^T$, so that the updated weights are $W_t = W_{0,t} + \Delta W_t$. We denote the resulting adapted network as $f_W$. 

    We consider the multi-label classification problem with an unknown data distribution $\mathcal{D}$ and a set of samples $S = \{(x_1, y_1), \dots, (x_N, y_N)\}\sim \mathcal{D}$, such that $\|x_i\|_2 \leqslant B_x$ and $y_i \in \{1, \dots, K\}$. For an input $x$, we denote the intermediate forward-pass states as
    $h_0 = x$ and $h_t = \sigma(W_t h_{t-1})$, with the final network output being $f_W(x) = h_T$.
    
    Additionally, we require limited spectral norms for the network before and after the adaptation: \\
    $\|W_{0,t}\|_{op} \leqslant B_t$, $\|\Delta W_t\|_{op} \leqslant \Lambda_t$, and $\|W_t\|_{op} \leqslant \bar{\beta}_t = B_t + \Lambda_t$. We assume that the stable rank of the adapters is bounded by a constant $\sr(\Delta W_t) \leqslant S_t$. The maximum layer width across the network is denoted as $d_{max} = \max_t (d_t)$, and the maximum number of parameters in any layer is $D_{max} = \max_t (d_t \cdot d_{t-1})$.\\

    For the main result in this section we use the \textit{Margin Loss}. First, we define the margin $m$ for the model $f_W$ and a data point $(x, y)$ as follows:
    $$ m(f_W(x), y) = f_W(x)_y - \max_{j \neq y} f_W(x)_j $$
    Then the true $\gamma$-margin error for $f_W$ is defined as:
    $$L_\gamma(W) = \mathbb{P}_{(x,y) \sim \mathcal{D}}\left[ m(f_W(x), y) \leqslant \gamma \right]$$
    Here, $L_0(W)$ represents the expected incorrect classification probability for the network (the true risk function). \\
    The empirical $\gamma$-margin error over the training set is:
    $$\widehat{L}_\gamma(W) = \frac{1}{N} \sum_{i=1}^N \mathbb{I}\left[ m(f_W(x_i), y_i) \leqslant \gamma \right]$$



    %%%%%%%%%%%%%%%%%%%%%%%%%%%%%%%%%%%%%%%%%%%%%%%%%%%%%%%%%%%%%%%%%%%%%%%%%%%%%%%%
    \subsubsection{Upper-bound for LoRA Generalization Gap}

    The following theorem formally derives the generalization bound incorporating the layer-wise stable rank constants $S_t$. The result isolates exactly how reducing stable rank acts as a powerful regularizer, naturally minimizing the generalization gap independently of the nominal dimensions.

    \begin{theorem}$~$

        For any $\gamma > 0$ and any $W = W_0 + \Delta W$ that satisfies the aforementioned conditions with probability at least $1-\delta$ by the choice of training sample S. The true classification risk is bounded by the following:
        $$ L_0(W) \le \hat{L}_\gamma(W) + \mathcal{O} \left( \sqrt{\frac{ B_x^2 T^2 (d_{max} + \ln(TN)) \left( \prod_{t=1}^T \bar{\beta}_t^2 \right) \sum_{l=1}^L \frac{\Lambda_l^2 S_t}{\bar{\beta}_t^2} + \log\left(\frac{N}{\delta}\right) }{\gamma^2 N}} \right) $$

    \end{theorem}
    
    \begin{proof}$~$

        We start with the fundamental PAC-Bayes theorem \cite{mcallester1999pac}\cite{mcallester2003simplified}:
        
        For any prior predictor distribution $\mathcal{P}$ (not depending on $S$) and any $\delta \in (0,1)$, with probability at least $1-\delta$, for any predictor distribution $\mathcal{Q}$ (depending on $S$) and the margin $m = \gamma/2$ error function the following holds:
        $$ \mathbb{E}_{w \sim \mathcal{Q}}[L_{\gamma/2}(w)] \leqslant \mathbb{E}_{w \sim \mathcal{Q}}[\hat{L}_{\gamma/2}(w)] + \sqrt{\frac{KL(\mathcal{Q} || \mathcal{P}) + \log\left(\frac{2N}{\delta}\right)}{2N}} $$

        Now we need to choose the distributions to be able to remove the expectations $\mathbb{E}_{w \sim \mathcal{Q}}$.\\
        
        Let $\mathcal{Q}$ be a distribution such that layer $t$ weights matrix is $w_t = W_t + U_t$ and $U_t\sim \mathcal{N}(0, \tau_t^2 I)$. The exact value of variance $\tau_t^2$ is to be determined.

        For given noise sample $U$ consider the maximum output perturbation over the dataset:
        $$ \Delta f(U) = \max_{x \in S} \|f_{W+U}(x) - f_W(x)\|_\infty $$

        Assume $\Delta f(U) \le \gamma / 4$, then for any $(x, y)$:
        $$|m(f_W(x), y) - m(f_{W+U}(x), y)| \leqslant \gamma/2$$
        Thus, if $f_W$ makes mistake $m(f_W(x), y) \leqslant 0$, then $m(f_{W+U}(x), y)\leqslant \gamma / 2$ which results in the following risks inequality:
        $$L_0(W) \leqslant L_{\gamma/2}(W+U)$$
        Similarly, the following holds:
        $$\widehat{L}_{\gamma/2}(W+U) \leqslant \widehat{L}_\gamma(W)$$
        
        So, combining with the assumption we get 
        $$\mathbb{I}[m(f_W(x), y)\leqslant 0] \leqslant \mathbb{I}[m(f_{W+U}(x), y)\leqslant \gamma/2] + \mathbb{I}[\Delta f(U) > \gamma / 4]$$
        and after taking the expectation:
        $$L_0(W)\leqslant \mathbb{E}_U[L_{\gamma / 2}(W + U)] + \mathbb{P}_U[\Delta f(U) > \gamma / 4]$$
        by doing the same, but averaging over $S$ instead of taking the expectation we get:
        $$\mathbb{I}[m(f_{W+U}(x), y)\leqslant \gamma / 2] \leqslant \mathbb{I}[m(f_{W}(x), y)\leqslant \gamma] + \mathbb{I}[\Delta f(U) > \gamma / 4]$$
        $$\mathbb{E}_U[\widehat{L}_{\gamma / 2}(W + U)] \leqslant \widehat{L}_\gamma(W) + \mathbb{P}_U[\Delta f(U) > \gamma / 4]$$

        Here we substitute these inequalities into the fundamental theorem:
        $$L_0(W)\leqslant \mathbb{E}_U[L_{\gamma / 2}(W + U)] + \mathbb{P}_U[\Delta f(U) > \gamma / 4]\leqslant \mathbb{E}_U[\widehat{L}_{\gamma / 2}(W + U)] + \mathbb{P}_U[\Delta f(U) > \gamma / 4] + \sqrt{\frac{KL(\mathcal{Q} || \mathcal{P}) + \log\left(\frac{2N}{\delta}\right)}{2N}}  $$
        $$L_0(W)\leqslant \widehat{L}_\gamma(W) + 2\mathbb{P}_U[\Delta f(U) > \gamma / 4] + \sqrt{\frac{KL(\mathcal{Q} || \mathcal{P}) + \log\left(\frac{2N}{\delta}\right)}{2N}}$$


        Now we must establish a strict upper bound on the perturbation $\Delta f(U)$ to determine the appropriate variance $\tau_t^2$ that makes $\mathbb{P}_U[\Delta f(U) > \gamma / 4]$ negligibly small.                

        \textbf{Bounding the network perturbation:}

        Let $h_t(x)$ denote the activations of the deterministic network $f_W$ at layer $t$, and $\tilde{h}_t(x)$ denote the activations of the noisy network $f_{W+U}$, let $h_0(x) = \tilde{h}_0(x) = x$. 
        
        For any subsequent layer $t$:
        $$ \|\tilde{h}_t - h_t\|_2 = \|\sigma((W_t + U_t)\tilde{h}_{t-1}) - \sigma(W_t h_{t-1})\|_2 $$

        Using the 1-Lipschitz property of the activation function $\sigma$ and the triangle inequality, we get:
        $$ \|\tilde{h}_t - h_t\|_2 \leqslant \|(W_t + U_t)\tilde{h}_{t-1} - W_t h_{t-1}\|_2 \leqslant \|W_t\|_{op} \|\tilde{h}_{t-1} - h_{t-1}\|_2 + \|U_t\|_{op} \|\tilde{h}_{t-1}\|_2 $$

        Since $\|W_t\|_{op} \leqslant \bar{\beta}_t$, we obtain a recursive inequality:
        $$ \|\tilde{h}_t - h_t\|_2 \leqslant \bar{\beta}_t \|\tilde{h}_{t-1} - h_{t-1}\|_2 + \|U_t\|_{op} \|\tilde{h}_{t-1}\|_2 $$

        Unrolling this recursion from layer $T$ down to layer 1, we obtain the bound for the final network output:
        $$ \|f_{W+U}(x) - f_W(x)\|_2 = \|\tilde{h}_T - h_T\|_2 \leqslant \sum_{t=1}^T \left( \prod_{i=t+1}^T \bar{\beta}_i \right) \|U_t\|_{op} \|\tilde{h}_{t-1}\|_2 $$

        To bound the norm of the noisy activations $\|\tilde{h}_{t-1}\|_2$, we do the following:
        $$ \|\tilde{h}_{t-1}\|_2 \leqslant \|x\|_2 \prod_{i=1}^{t-1} \|W_i + U_i\|_{op} \leqslant B_x \prod_{i=1}^{t-1} (\bar{\beta}_i + \|U_i\|_{op}) $$

        Assuming that we can choose the noise variance small enough (we do that below) such that for all layers $\|U_i\|_{op} \leqslant \frac{1}{T} \bar{\beta}_i$, we can bound the product:
        $$ \bar{\beta}_i + \|U_i\|_{op} \leqslant \bar{\beta}_i \left(1 + \frac{1}{T}\right) $$

        $$ \|\tilde{h}_{t-1}\|_2 \leqslant B_x \left(1 + \frac{1}{T}\right)^{t-1} \prod_{i=1}^{t-1} \bar{\beta}_i \leqslant e B_x \prod_{i=1}^{t-1} \bar{\beta}_i $$

        Substituting this back into the unrolled recursion, we get:
        $$ \|f_{W+U}(x) - f_W(x)\|_2 \leqslant \sum_{t=1}^T \left( \prod_{i=t+1}^T \bar{\beta}_i \right) \|U_t\|_{op} \left( e B_x \prod_{i=1}^{t-1} \bar{\beta}_i \right) = e B_x \left( \prod_{i=1}^T \bar{\beta}_i \right) \sum_{t=1}^T \frac{\|U_t\|_{op}}{\bar{\beta}_t} $$

        Since we bounded $\|x\|_2$ with $B_x$ we have already taken the maximum over $x \in S$, so the final perturbation bound $\Delta f(U)$:
        $$ \Delta f(U) \leqslant e B_x \left( \prod_{i=1}^T \bar{\beta}_i \right) \sum_{t=1}^T \frac{\|U_t\|_{op}}{\bar{\beta}_t} $$


        \textbf{Bounding the spectral norm of U:} 

        Since the noise matrix $U_t \in \mathbb{R}^{d_t \times d_{t-1}}$ has entries from $\mathcal{N}(0, \tau_t^2)$ we can apply Gordon's theorem for gaussian matrices\cite{vershynin2018high}, for any $s > 0$, with probability at least $1 - \exp(-s^2/2)$:
        $$ \|U_t\|_{op} \leqslant \tau_t (\sqrt{d_t} + \sqrt{d_{t-1}} + s) \leqslant \tau_t (2\sqrt{d_{max}} + s) $$
        For all $T$ layers simultenously this hold with probability at least $1 - T\exp(-s^2/2)$.\\


        \textbf{Chosing the variance and bounding the failure probability:}

        As mentioned above, we require the total perturbation to satisfy $\Delta f(U) \leqslant \gamma / 4$. Substituting the spectral norm bound into the perturbation inequality yields:
        $$ \Delta f(U) \leqslant e B_x \left( \prod_{i=1}^T \bar{\beta}_i \right) \sum_{t=1}^T \frac{\tau_t (2\sqrt{d_{max}} + s)}{\bar{\beta}_t} $$

        To balance the noise across layers, we set the variance multiplier proportional to the spectral norm bound of each layer: $\tau_t = \tau \bar{\beta}_t$. The sum then simplifies:
        $$ \sum_{t=1}^T \frac{\tau \bar{\beta}_t (2\sqrt{d_{max}} + s)}{\bar{\beta}_t} = \tau T (2\sqrt{d_{max}} + s) $$

        Thus, the perturbation is bounded by:
        $$ \Delta f(U) \leqslant e B_x \left( \prod_{i=1}^T \bar{\beta}_i \right) \tau T (2\sqrt{d_{max}} + s) $$

        We set this expression exactly equal to $\gamma / 4$ and solve for the global variance data dependent parameter $\tau^*$:
        $$ \tau^* = \frac{\gamma}{4 e B_x T (2\sqrt{d_{max}} + s) \prod_{i=1}^T \bar{\beta}_i} $$

        Now, we additionally require the probability of the perturbation exceeding $\gamma/4$ to be negligibly small, specifically 
        $$\mathbb{P}_U[\Delta f(U) > \gamma/4] \leqslant \frac{1}{\sqrt{N}}$$ 
        
        Setting the failure probability from the $\|U\|_{op}$ bound equal to this value to find the specific $s$:
        $$ T\exp(-s^2/2) = \frac{1}{\sqrt{N}} \implies s = \sqrt{2 \ln(T \sqrt{N})} $$


        \textbf{Calculating KL Divergence and Introducing Stable Rank:} 

        Let the distribution $\mathcal{P}$ be the $\prod_{t=1}^T \mathcal{N}(W_{0,t}, \tau_t^2 I)$. Distribution $\mathcal{Q}$ as previously mention is $\prod_{t=1}^T \mathcal{N}(W_{0,t} + \Delta W_t, \tau_t^2 I)$. The KL divergence between these isotropic multivariate gaussians is the following:
        $$ KL(\mathcal{Q} || \mathcal{P}) = \sum_{t=1}^T \frac{\|(W_{0,t} + \Delta W_t) - W_{0,t}\|_F^2}{2\tau_t^2} = \sum_{t=1}^T \frac{\|\Delta W_t\|_F^2}{2\tau_t^2} $$

        Substituting $\tau_t^2 = \tau^2 \bar{\beta}_t^2$, we obtain:
        $$ KL(\mathcal{Q} || \mathcal{P}) = \frac{1}{2\tau^2} \sum_{t=1}^T \frac{\|\Delta W_t\|_F^2}{\bar{\beta}_t^2} $$

        Here we finally utilize the definition of stable rank: $\sr(\Delta W_t) = \frac{\|\Delta W_t\|_F^2}{\|\Delta W_t\|_{op}^2} \leqslant S_t$. Since the spectral norm of the adapter is bounded by $\|\Delta W_t\|_{op} \leqslant \Lambda_t$, then the Frobenius norm is bounded by:
        $$ \|\Delta W_t\|_F^2 \leqslant \Lambda_t^2 S_t $$
        Substituting this inequality into the KL divergence yields:
        $$ KL(\mathcal{Q} || \mathcal{P}) \leqslant \frac{1}{2\tau^2} \sum_{t=1}^T \frac{\Lambda_t^2 S_t}{\bar{\beta}_t^2} $$


        \textbf{The union bound over $\tau$ (PAC-Bayes prior fixing):} 

        A fundamental requirement of the PAC-Bayes framework is that the prior distribution $\mathcal{P}$ (and consequently $\tau$) must be chosen \textit{before} observing the training data $S$. However, the optimal variance $\tau^*$ depends on the spectral norms $\bar{\beta}_t$ and the actual learned parameters, making it data-dependent. We cannot directly use $\tau^*$ to define a single prior without violating the PAC-Bayes independence assumption.
        
        To resolve this rigorously, we employ a standard union bound technique over a predefined grid of priors. Before observing the data, we define a countably infinite grid of variance multipliers:
        $$ \tau_{(k)} = \exp(k/T), \quad \text{for } k \in \mathbb{Z} $$
        For each $k$, we construct a specific prior $\mathcal{P}_k$ characterized by the layer-wise variances $\tau_{t,(k)}^2 = \tau_{(k)}^2 \bar{\beta}_t^2$.
        
        We must distribute our total allowed failure probability $\delta$ across all these priors. We assign a confidence level $\delta_k$ to the $k$-th prior such that:
        $$ \delta_k = \frac{\delta}{2(|k|+1)^2},~~~  \sum_{k=-\infty}^{\infty} \delta_k \leqslant \delta$$

        By applying the union bound, the fundamental PAC-Bayes inequality (established at the beginning of the proof) holds \textit{simultaneously} for all priors $\mathcal{P}_k \in \{\mathcal{P}_k\}_{k \in \mathbb{Z}}$ with probability at least $1-\delta$.
        
        After the training process is complete and the adapters $\Delta W_t$ are obtained, we calculate the theoretically optimal data-dependent variance $\tau^*$. We then simply select the integer index $k$ such that our grid value closely approximates the optimal value:
        $$ \tau_{(k)} \leqslant \tau^* < \tau_{(k+1)} $$
        Because $\tau_{(k+1)} / \tau_{(k)} = \exp(1/T)$, it is guaranteed that $\tau_{(k)} > \tau^* e^{-1/T}$. 

        Substituing this into $KL$-divergence considering $T\geqslant 1 \implies e^{2/T} \leqslant e^2$, yields the following bound:
        $$ KL(\mathcal{Q} || \mathcal{P}_k) = \frac{1}{2\tau_{(k)}^2} \sum_{t=1}^T \frac{\|\Delta W_t\|_F^2}{\bar{\beta}_t^2} \leqslant \frac{e^2}{2(\tau^*)^2} \sum_{t=1}^T \frac{\Lambda_t^2 S_t}{\bar{\beta}_t^2} $$

        Furthermore, in the PAC-Bayes bound, the logarithmic term now uses $\delta_k$ instead of $\delta$:
        $$ \log\left(\frac{2N}{\delta_k}\right) = \log\left(\frac{2N}{\delta}\right) + \log(2(|k|+1)^2) $$
        To understand the magnitude of this penalty, recall that we chose $k$ such that $\exp(k/T) \approx \tau^*$. Therefore, the index magnitude $|k|$ is bounded by $\mathcal{O}(T |\ln(\tau^*)|)$. Looking at the closed-form expression for $\tau^*$, $|\ln(\tau^*)|$ scales logarithmically with the network dimensions and the product of spectral norms $|\ln(\tau^*)| = \mathcal{O}(\ln(d_{max}) + \sum_{i=1}^T \ln(\bar{\beta}_i))$. 
        
        Consequently, the penalty term $\log(2(|k|+1)^2)$ introduces a double-logarithmic dependence on the network parameters and norms:
        $$ \log(2(|k|+1)^2) = \mathcal{O} \left( \log\left( T \ln(d_{max}) + T \sum_{i=1}^T \ln(\bar{\beta}_i) \right) \right) $$
        This $\log\log(\cdot)$ penalty grows so slowly that it behaves as a negligible $\tilde{\mathcal{O}}(1)$ constant for any practically realizable neural network.


        \textbf{Final substitution and asymptotic bound:} 

        Now we substitute the expression for $\tau^*$ into our bounded KL divergence:
        $$ \frac{1}{(\tau^*)^2} = \frac{16 e^2 B_x^2 T^2 \left(2\sqrt{d_{max}} + \sqrt{2 \ln(T \sqrt{N})}\right)^2 \prod_{t=1}^T \bar{\beta}_t^2}{\gamma^2} $$
        Using the basic algebraic inequality $(a+b)^2 \leqslant 2a^2 + 2b^2$, we obtain the following:
        $$ \left(2\sqrt{d_{max}} + \sqrt{2 \ln(T \sqrt{N})}\right)^2 \leqslant 8 d_{max} + 4 \ln(T \sqrt{N}) $$
        As $N \geqslant 1$, we can derive another bound:
        $$\ln(T \sqrt{N}) = \ln(T) + \frac{1}{2}\ln(N) \leqslant \ln(TN)$$
        Thus, the KL divergence is then asymptotically bounded by:
        $$ KL(\mathcal{Q} || \mathcal{P}_k) = \mathcal{O} \left( \frac{B_x^2 T^2 (d_{max} + \ln(TN)) \left( \prod_{t=1}^T \bar{\beta}_t^2 \right)}{\gamma^2} \sum_{t=1}^T \frac{\Lambda_t^2 S_t}{\bar{\beta}_t^2} \right) $$
        Finally, we substitute the KL bound and our previously fixed $\mathbb{P}_U \leqslant 1/\sqrt{N}$ back into the fundamental PAC-Bayes inequality established at the beginning of the proof. Since the term $2/\sqrt{N} = \mathcal{O}(1/\sqrt{N})$ scales identically to the denominator of the square root, it is cleanly absorbed into the main $\mathcal{O}(\cdot)$ notation. Taking the square root of the KL divergence expression yields the final bound:
        $$ L_0(W) \leqslant \hat{L}_\gamma(W) + \mathcal{O} \left( \sqrt{\frac{ B_x^2 T^2 (d_{max} + \ln(TN)) \left( \prod_{t=1}^T \bar{\beta}_t^2 \right) \sum_{t=1}^T \frac{\Lambda_t^2 S_t}{\bar{\beta}_t^2} + \log\left(\frac{N}{\delta}\right) }{\gamma^2 N}} \right) $$
        This completes the proof.

    \end{proof}

    While the derived PAC-Bayesian bound remains purely theoretical and potentially quite loose compared to practical applications -- as it aggregates worst-case layer-wise perturbations without fully accounting for complex cross-layer interactions -- Theorem 3 provides a valuable theoretical perspective on the generalization dynamics of LoRA. The bound suggests a mathematical rationale for why tuning large nominal ranks not necessarily leads to catastrophic overfitting in practice. Within this framework, the generalization gap correlates fundamentally with the stable rank sum rather than the algebraic rank $r$. If the training dynamics implicitly bias the adapter updates toward highly skewed singular value distributions ($\sr \ll \rank $), the effective capacity of the network remains tightly constrained regardless of the parameter budget. Consequently, this stable-rank dependence offers a theoretical hypothesis for how LoRA resists overfitting, suggesting that monitoring stable rank could serve as a useful diagnostic tool for predicting and analyzing generalization performance across different optimization techniques or architectural modifications.

    Before concluding this section, we must acknowledge a slight formal distinction between our PAC-Bayesian analysis and the framework established by Kratsios et al. \cite{kratsios2025sharpgeneralizationboundsfoundation}. While their uniform convergence bound explicitly targets the standard generalization gap (the absolute difference between empirical and expected risk) for a bounded 1-Lipschitz loss, our Theorem 3 bounds the true expected classification error $L_0(W)$ using an empirical $\gamma$-margin loss $\widehat{L}_\gamma(W)$. Because these bounds govern slightly different mathematical formulations of the learning objective, they are not strictly one-to-one comparable. 
    However, within the margin-based framework, the margin parameter $\gamma$ acts as a tunable strictness threshold; by adjusting $\gamma$ to be arbitrarily small, the margin setup can conceptually approximate the standard continuous loss behavior with the required scaling. Most importantly, this shift in the loss definition does not alter the primary comparative takeaway of our analysis: whereas the covering-number approach permanently locks the complexity penalty to the algebraic rank $\mathcal{O}(\sqrt{r/N})$, the PAC-Bayesian margin framework successfully transitions the scaling to the stable rank, proving that the true generalization penalty is governed by the spectral distribution of the adapter rather than its parameter count.



    %%%%%%%%%%%%%%%%%%%%%%%%%%%%%%%%%%%%%%%%%%%%%%%%%%%%%%%%%%%%%%%%%%%%%%%%%%%%%%%%
    \section{Empirical Validation}
    % ==============================================================================
    % IN SHORT:
    % - Toy experiments on synthetic data validate the main theoretical claim: generalization gap scales with √SR, not √rank.
    % - 10-layer MLPs in a teacher-student setup; LoRA adapters trained under controlled target shift complexity.
    % - Regression (vector output, MSE) is cleaner and more stable; classification (0-1 error, CE loss) mirrors Theorem 3 setup.
    % - Both experiments confirm trained adapter SR saturates far below nominal rank, and generalization gap aligns with √SR across conditions.
    % - Full fine-tuning baseline consistently shows higher SR and larger generalization gap.
    % ==============================================================================
    
    To empirically validate our theoretical claims---in particular, the generalization bound of Theorem~3 which predicts that the generalization gap scales as $\mathcal{O}(\sqrt{\sr(\Delta W)/N})$ rather than $\mathcal{O}(\sqrt{\rank(\Delta W)/N})$ -- we design controlled synthetic experiments where the complexity of the target adaptation is systematically varied. The core objective is to demonstrate that, across a range of target difficulties, the observed generalization gap of LoRA adapters is practically governed by the trained stable rank, not the nominal adapter rank. We additionally confirm the implicit low-SR bias of LoRA optimization, showing that even when the nominal rank is increased up to the full layer width, the trained stable rank saturates at a value substantially smaller than $r$.


    %%%%%%%%%%%%%%%%%%%%%%%%%%%%%%%%%%%%%%%%%%%%%%%%%%%%%%%%%%%%%%%%%%%%%%%%%%%%%%%%
    \subsection{Experimental Setup}

    \subsubsection*{Network architecture and teacher-student framework}

    We employ a teacher-student setup in which both models share the same MLP architecture of width $D=128$ with $L=10$ linear layers and SiLU activations. The teacher is constructed by taking the (random) student base and manually shifting a subset of its weight layers with a controlled perturbation $\delta_l$. The student is then trained to recover the teacher's outputs on a finite sample, using LoRA adapters injected into those same perturbed layers. Crucially, this is not knowledge distillation (which would involve matching logit distributions via KL divergence); rather, the teacher simply serves as a data-generating process: its inputs and noisy outputs are used as $(x, y)$ training pairs for the student. 

    We deliberately choose a deep 10-layer architecture because preliminary experiments showed that shallower networks do not exhibit the low-SR implicit bias as prominently. That might be because the deeper model better approximates the optimization dynamics of practically used large networks such as LLMs, where the low-SR phenomenon has been documented~\cite{lion2025polar}.

    \subsubsection*{Controlling target complexity via stable rank}

    The key experimental variable is the stable rank of the target weight shift $\delta_l$ applied to each perturbed layer. We generate shifts via a controlled SVD procedure: random orthogonal matrices $U$, $V$ are drawn and $\delta = U \operatorname{diag}(\sigma) V^T$ is constructed with singular values set so that $\sigma_1$ equals a chosen \textbf{shift multiplier} and $\sigma_{i \geqslant 2} \approx 1$, after which the singular value vector is normalized to unit norm. We use eight conditions $c_1$ through $c_8$ corresponding to shift multipliers $\{1.0,\;1.5,\;2.0,\;3.0,\;5.0,\;8.0,\;12.0,\;20.0\}$. A multiplier of $1.0$ produces a shift with stable rank close to the layer width $D$ (all singular values nearly equal), while a multiplier of $20.0$ concentrates virtually all energy in the first singular component, yielding $\sr \approx 1$. Each shift is additionally scaled by a factor ($0.5$ for regression, $0.03$ for classification) and proportionally to the RMS of the base layer weights. This design gives us direct control over the intrinsic complexity of the training signal: a high-SR shift is harder for a student to match because the student's optimizer is implicitly biased toward low-SR solutions.

    \subsubsection*{Training protocol and metrics}

    All experiments use $N_{\text{train}} = 2048$ training samples and $N_{\text{val}} = 10240$ validation samples drawn from a standard normal distribution. Gaussian noise is added to teacher outputs (standard deviation $\sigma_{\text{noise}} = 0.15$ for regression, $0.01$ for classification) to induce overfitting and add some compelixy and diversity. LoRA adapters are trained at nominal ranks $r \in \{4, 6, 8, 12, 16, 24, 32, 48, 64, 96, 128\}$, covering the full range up to the layer width. A full fine-tuning (Full FT) baseline is also included, where the same subset of layers is trained without rank constraints. All models are optimized with AdamW for 200 epochs with a learning rate of $10^{-3}$.

    The primary metric is the \textbf{generalization gap}, defined as the difference between validation and training loss. For each trained adapter, we also record the \textbf{average stable rank} $\overline{\sr} = \frac{1}{|\mathcal{L}_{\text{adpt}}|} \sum_{l \in \mathcal{L}_{\text{adpt}}} \sr(\Delta W_l)$ computed directly from the learned $\Delta W_l = B_l A_l$ matrices.

    We conduct two variants of the experiment. The first is a \textbf{regression} task: the teacher outputs a $D$-dimensional real vector, and the student minimizes MSE. The second is a \textbf{classification} task: the teacher produces logits over $K=64$ classes, hard labels are obtained via argmax, and the student is trained with cross-entropy loss; generalization is measured as the 0-1 error gap. Although the regression setup does not replicate the exact formulation of Theorem~3, it is similar enough while being substantially more stable numerically, making it the cleaner demonstration of the core idea. The classification experiment more directly mirrors the theoretical setting of the PAC-Bayesian margin bound. (Runs shown: 2048 training samples with Full FT baseline for regression; 10-layer, 64-class for classification.)


    %%%%%%%%%%%%%%%%%%%%%%%%%%%%%%%%%%%%%%%%%%%%%%%%%%%%%%%%%%%%%%%%%%%%%%%%%%%%%%%%
    \subsection{Regression Experiment}

    In the regression variant, the teacher and student both output $D$-dimensional real vectors, with shifts applied to 9 intermediate layers (layers 1 through 9) and LoRA adapters injected into the same layers. Noise with $\sigma_{\text{noise}}=0.15$ is added to the teacher outputs.

    \begin{wrapfigure}{l}{0.5\textwidth}
        \centering
        \includegraphics[width=0.48\textwidth]{"../figures/regression_trained_sr_vs_nominal_rank.png"}
        \caption{Trained stable rank versus nominal LoRA rank across target shift conditions.}
        \label{fig:regr_sr_vs_rank}
    \end{wrapfigure}

    Figure~\ref{fig:regr_sr_vs_rank} shows that the trained adapter stable rank $\overline{\sr}$ increases with nominal rank $r$ but saturates early -- around $r \approx 32$--$48$ -- and remains nearly flat thereafter, even when $r$ reaches the full layer width of 128. However, saturation alone does not establish an implicit low-SR bias -- it could simply reflect interlayer communication compensating for low per-layer SR. The decisive evidence for the implicit bias is the contrast with Full FT: Full fine-tuning of the network, even though having the same algebraic rank, achieves much higher stable rank and lower training loss, that incurs a larger generalization gap (overfitting), whereas LoRA adapters remain at lower SR with a smaller gap, indicating that the optimization dynamics bound effective capacity in a manner that favors generalization. Notably, the SR curves for all eight conditions are closely clustered together, indicating that the trained SR in this particular scenario is primarily determined by the nominal rank and the optimization process, with the target shift's SR exerting limited influence on the adapter's saturation level. The Full FT baselines (dashed lines) consistently achieve substantially higher stable rank ($\overline{\sr} \approx 27$--$29$) than any LoRA configuration, with also negligible variation across conditions. \\

    \begin{figure}[htbp]
        \centering
        \begin{minipage}[t]{0.5\textwidth}
            \centering
            \includegraphics[width=\textwidth]{"../figures/regression_gen_gap_vs_sqrt_trained_sr.png"}
            \caption{Generalization gap versus square root of trained stable rank.}
            \label{fig:regr_gap_vs_sr}
        \end{minipage}%
        \begin{minipage}[t]{0.5\textwidth}
            \centering
            \includegraphics[width=\textwidth]{"../figures/regression_gen_gap_vs_sqrt_rank.png"}
            \caption{Generalization gap versus square root of nominal rank.}
            \label{fig:regr_gap_vs_rank}
        \end{minipage}
    \end{figure}

    Figure~\ref{fig:regr_gap_vs_sr} plots the generalization gap (validation MSE $-$ training MSE) against $\sqrt{\overline{\sr}}$. The data points from all conditions and all nominal ranks follow a broadly monotonic trend: higher trained SR is associated with a larger generalization gap. The Full FT points (star markers) lie at the high-SR, high-gap end of the distribution, consistent with their larger effective capacity. This trend is the central empirical finding: the generalization gap is primarily a function of the trained stable rank, not of the nominal parameter budget. \\
    Figure~\ref{fig:regr_gap_vs_rank} provides the contrast. When the generalization gap is plotted against $\sqrt{r}$, the curves remain closely packed across all conditions -- there is still no separation by condition. Unlike Figure~\ref{fig:regr_gap_vs_sr} where the gap grows with $\sqrt{\overline{\sr}}$, here plotted against $\sqrt{r}$ the gap plateaus: increasing nominal rank does not systematically increase the generalization gap. This reinforces the conclusion from Figure~\ref{fig:regr_gap_vs_sr} that the generalization gap is tied to $\sqrt{\sr}$, not $\sqrt{r}$.
    

    %%%%%%%%%%%%%%%%%%%%%%%%%%%%%%%%%%%%%%%%%%%%%%%%%%%%%%%%%%%%%%%%%%%%%%%%%%%%%%%%
    \subsection{Classification Experiment}

    The classification variant aligns more closely with the margin-based PAC-Bayesian setup of Theorem~3. The teacher outputs logits over $K=64$ classes; hard labels are obtained by taking argmax of noise-corrupted logits ($\sigma_{\text{noise}}=0.01$). LoRA adapters and shifts are applied to 3 intermediate layers. The generalization gap is measured as the 0-1 error difference (validation error $-$ training error).

    \begin{wrapfigure}{l}{0.5\textwidth}
        \centering
        \includegraphics[width=0.5\textwidth]{"../figures/classification_trained_sr_vs_nominal_rank.png"}
        \caption{Trained stable rank versus nominal LoRA rank across target shift conditions.}
        \label{fig:cls_sr_vs_rank}
    \end{wrapfigure}

    Figures~\ref{fig:cls_sr_vs_rank}--\ref{fig:cls_gap_vs_rank} show the corresponding results. The classification setting is inherently noisier than regression---a consequence of the discrete 0-1 error surface and the indirect relationship between cross-entropy optimization and the classification margin---yet the same qualitative patterns emerge. Figure~\ref{fig:cls_sr_vs_rank} confirms that the trained SR saturates at values far below the nominal rank (in the range of ~1.4--3.5 across all conditions and ranks). Figure~\ref{fig:cls_gap_vs_sr} reveals that, despite higher variance, the generalization gap exhibits a discernible positive association with $\sqrt{\overline{\sr}}$. Figure~\ref{fig:cls_gap_vs_rank} shows that $\sqrt{r}$ does not provide a consistent predictor and the overall behaviour is closer to the plateau of figure~\ref{fig:regr_gap_vs_rank} than to the noisy growth of generalisation gap vs SR in figure~\ref{fig:cls_gap_vs_sr}.
    \begin{figure}[htbp]
        \centering
        \begin{minipage}[t]{0.5\textwidth}
            \centering
            \includegraphics[width=\textwidth]{"../figures/classification_gen_gap_vs_sqrt_trained_sr.png"}
            \caption{Generalization gap versus square root of trained stable rank.}
            \label{fig:cls_gap_vs_sr}
        \end{minipage}%
        \begin{minipage}[t]{0.5\textwidth}
            \centering
            \includegraphics[width=\textwidth]{"../figures/classification_gen_gap_vs_sqrt_rank.png"}
            \caption{Generalization gap versus square root of nominal rank.}
            \label{fig:cls_gap_vs_rank}
        \end{minipage}
    \end{figure}

    Notably, the observed 0-1 error differences are modest in absolute terms, and the Full FT baseline achieves near-zero training error with a large gap to validation error ($\approx 0.73$--$0.77$), while some of the LoRA adapters achieve substantially smaller gaps. The practical takeaway remains: \textbf{the generalization behavior is better characterized by $\sqrt{\sr(\Delta W)}$ than by $\sqrt{\rank(\Delta W)}$, as predicted by Theorem~3.} The noise in the classification curves underscores that, in settings with discrete losses, the margin-based bound captures a trend rather than an exact functional relationship, and larger-scale experiments would likely be needed for a tighter quantitative verification. \\



    %%%%%%%%%%%%%%%%%%%%%%%%%%%%%%%%%%%%%%%%%%%%%%%%%%%%%%%%%%%%%%%%%%%%%%%%%%%%%%%%
    \subsection{Analysis and Discussion}
    \label{sec:analysis}

    Across both regression and classification, the generalization gap is consistently better predicted by $\sqrt{\overline{\sr}}$. The optimizer's implicit bias toward low-SR solutions is robust to both task type and target complexity: even when the data is generated with high-SR shifts and expected to require higher SR solution, LoRA tends to converge to low-SR ones, even though full fine-tuning consistently produces higher SR updates. Thus, rank thus systematically overestimates LoRA's effective capacity.

    As a future research direction, we would add that beyond post-hoc diagnostics, the SR saturation curve can guide rank selection upfront: by measuring where SR plateaus for a given architecture, practitioners can identify the minimal effective rank without sweeping validation performance across ranks. Additionally, SR provides a unified metric for comparing LoRA variants or optimization techniques -- a higher SR at the same nominal rank apparently indicates better utilization of the allocated parameter budget.

    Extending these measurements to LLM-scale LoRA fine-tuning on NLP benchmarks would test whether the SR--generalization relationship holds at scale. Investigating optimizer-specific effects on the SR distribution could reveal whether the low-SR bias is optimizer-dependent or a deeper property of factored gradient descent.




    %%%%%%%%%%%%%%%%%%%%%%%%%%%%%%%%%%%%%%%%%%%%%%%%%%%%%%%%%%%%%%%%%%%%%%%%%%%%%%%%
    \section{Conclusion}

    % Open Questions: Why is SR low? Is it a fundamental property of gradient descent in overparameterized regimes? Does the optimizer (AdamW vs SGD vs Muon) dictate this bias or some other underlying mechanism for broad class of gradient-based optimization techniques? 
    % Note that AdamW's coordinate-wise scaling might inherently suppress non-dominant singular directions, whereas something like Muon (which applies orthogonal momentum) might yield a different SR distribution. (or not?)
    % Extendability: How does this apply to other PEFT methods...? Since they alter initialization but retain low-rank gradients, the SR framework should directly apply. 
    % Bla-bla other extensions/applications... 
    
    

    In this thesis, we established stable rank as the mathematically appropriate theoretical measure for analyzing LoRA's properties, correcting the persistent reliance on algebraic rank in existing literature. Our key contributions are:
    
    Fundamental Matrix Limits (Lemmas 1-2): We derived strict bounds on matrix approximation under stable rank constraints, demonstrating an irreducible Frobenius error when attempting to approximate high-SR targets with low-SR matrices. We also provied explicit optimal solution to the rank- and stable rank-constrained approximation problem (in terms of minimum Frobenius error).

    Expressivity Bounds (Theorems 1-2): using derived optimal solution for single-matrix approximation, we extend rank-based expressivity bounds \cite{zeng2024expressivepowerlowrankadaptation} for FNNs and transformers to stable rank-constrained scenario, that is widely common in practice. We claim that our bounds better represent the actual expressivity capacity scaling, even though the exact expression is still way too far from practical application.
    

    Generalization Bounds (Theorem 3): We leveraged a PAC-Bayesian framework to rigorously prove a tighter generalization gap that scales with $\mathcal{O}(\sqrt{\sum_t \Lambda_t^2 S_t / \bar{\beta}_t^2})$. Because optimization implicitly biases LoRA such that $S_t \ll r_t$, our theorem mathematically validates why massively overparameterized LoRA adapters generalize remarkably well in practice.
    
    % [TODO: Add final concluding sentences after the empirical section is written.]


    \bibliographystyle{plain}
    \bibliography{references.bib}

\end{document}
